\documentclass[aspectratio=169,8pt]{beamer}
\usetheme{Madrid}
\usepackage[utf8]{inputenc}
\usepackage[T1]{fontenc}
\usepackage{amsmath}
\usepackage{amssymb}
\usepackage{booktabs}
\usepackage{enumitem}

\title{GARP Risk and AI Exam Prep}
\subtitle{100 Practice Questions with Theory}
\author{Unofficial Study Aid}
\date{\today}

\begin{document}
	
	\begin{frame}
		\titlepage
	\end{frame}
	
	\begin{frame}{Outline}
		\tableofcontents
	\end{frame}
	
	% ============================================================
	\section{Module 1: AI and Risk}
	% ============================================================
	
	\begin{frame}{Q1 --- Inscrutability}
		\begin{block}{Question}
			A bank uses AI models for detecting fraudulent transactions. While the models are highly effective, the risk team is concerned about inscrutability. Which statement correctly describes inscrutability?
		\end{block}
		\begin{enumerate}[label=\Alph*.]
			\item Inability to understand the internal mechanisms of a model.
			\item Inability to determine the accuracy of a model's prediction.
			\item Inability to explain a model's predictions.
			\item Inability to generalize a model to out-of-sample data.
		\end{enumerate}
		\begin{exampleblock}{Answer: A}
		\end{exampleblock}
		\begin{block}{Explanation}
			Inscrutability = inability to understand internal mechanisms. Distinct from accuracy, explainability, and generalizability.
		\end{block}
	\end{frame}
	
	\begin{frame}{Q1 Theory --- Inscrutability vs Explainability vs Interpretability}
		\begin{block}{Theory}
			\textbf{Inscrutability} = cannot understand internal mechanisms.\\
			\textbf{Explainability} = can explain why a decision was made.\\
			\textbf{Interpretability} = can understand internal workings.
			\vspace{2mm}
			
			\textbf{Key points:}
			\begin{itemize}
				\item Deep neural nets are accurate but inscrutable.
				\item Inscrutability is a model risk (blocks validation, audit, compliance).
				\item GARP exam tests the difference between these three terms.
			\end{itemize}
		\end{block}
	\end{frame}
	
	\begin{frame}{Q2 --- Classical AI}
		\begin{block}{Question}
			Deep Blue defeated Kasparov in 1997. Which best describes a characteristic of classical AI?
		\end{block}
		\begin{enumerate}[label=\Alph*.]
			\item Classical AI optimizes using gradient descent.
			\item Classical AI cannot process time series data.
			\item Classical AI relies on pre-defined rules in simple contexts.
			\item Classical AI scales easily to complex problems.
		\end{enumerate}
		\begin{exampleblock}{Answer: C}
		\end{exampleblock}
		\begin{block}{Explanation}
			Classical AI (GOFAI) uses rules and logic. No gradient descent (that is deep learning). Does not scale to complex problems.
		\end{block}
	\end{frame}
	
	\begin{frame}{Q2 Theory --- Classical AI vs Deep Learning}
		\begin{block}{Theory}
			\textbf{Classical AI (GOFAI):}
			\begin{itemize}
				\item Rule-based, symbolic logic
				\item Search algorithms (A*, minimax, lookahead)
				\item Simple, well-defined contexts (chess, tic-tac-toe)
				\item Cannot scale to complex problems
			\end{itemize}
			\textbf{Deep Learning:}
			\begin{itemize}
				\item Neural networks, many layers
				\item Backpropagation + gradient descent
				\item Self-training (learns by competing against itself)
				\item Handles unstructured data (images, text, audio)
				\item Inscrutable (black-box)
			\end{itemize}
			\textbf{Exam tip:} ``Self-training'' or ``gradient descent'' = deep learning.
		\end{block}
	\end{frame}
	
	\begin{frame}{Q3 --- Societal Risk}
		\begin{block}{Question}
			Which best describes a societal risk scenario regarding AI?
		\end{block}
		\begin{enumerate}[label=\Alph*.]
			\item Trading company errors cause client losses.
			\item Social media company spreads misinformation influencing public opinion.
			\item Manufacturer hiring decisions are poor.
			\item Hospital accesses personal info causing privacy loss.
		\end{enumerate}
		\begin{exampleblock}{Answer: B}
		\end{exampleblock}
		\begin{block}{Explanation}
			Societal risk = harm to society at large. Misinformation influencing public opinion is societal.
		\end{block}
	\end{frame}
	
	\begin{frame}{Q3 Theory --- Individual, Organizational, Societal Risk}
		\begin{block}{Theory}
			\textbf{Three levels of AI risk:}
			\begin{itemize}
				\item \textbf{Individual:} harm to a person (privacy, wrong diagnosis).
				\item \textbf{Organizational:} harm to a company (losses, reputation, fines).
				\item \textbf{Societal:} harm to society (misinformation, inequality, trust erosion).
			\end{itemize}
			\textbf{Exam tip:} ``Society'', ``public opinion'', ``democracy'' = societal risk.
		\end{block}
	\end{frame}
	
	\begin{frame}{Q4 --- Deep Learning Feature}
		\begin{block}{Question}
			Which is a key differentiating feature of deep learning over classical AI?
		\end{block}
		\begin{enumerate}[label=\Alph*.]
			\item A* algorithm
			\item Recursive adversarial search
			\item Lookahead strategy
			\item Self-training
		\end{enumerate}
		\begin{exampleblock}{Answer: D}
		\end{exampleblock}
		\begin{block}{Explanation}
			Self-training is a deep learning feature. A*, adversarial search, and lookahead are classical AI.
		\end{block}
	\end{frame}
	
	\begin{frame}{Q4 Theory --- AI Timeline}
		\begin{block}{Theory}
			\textbf{AI breakthroughs:}
			\begin{itemize}
				\item 1950s--1980s: Classical AI (logic, search, expert systems)
				\item 1997: Deep Blue beats Kasparov
				\item 2010s: Deep learning revolution (AlexNet, AlphaGo)
				\item 2017: Transformers (``Attention Is All You Need'')
				\item 2020s: Large Language Models (GPT, Claude, Gemini)
			\end{itemize}
			\textbf{Deep learning features:}
			\begin{itemize}
				\item Self-training, backpropagation, gradient descent
				\item Automatic feature extraction
				\item Handles unstructured data
				\item Inscrutable (black-box)
			\end{itemize}
		\end{block}
	\end{frame}
	
	\begin{frame}{Q5 --- Data Types}
		\begin{block}{Question}
			A survey asks customers to select investment types (stocks, bonds, ETFs). How should these results be classified?
		\end{block}
		\begin{enumerate}[label=\Alph*.]
			\item Ordinal data
			\item Nominal data
			\item Interval data
			\item Ratio data
		\end{enumerate}
		\begin{exampleblock}{Answer: B}
		\end{exampleblock}
		\begin{block}{Explanation}
			Categories with no order = nominal data.
		\end{block}
	\end{frame}
	
	\begin{frame}{Q5 Theory --- Data Types}
		\begin{block}{Theory}
			\textbf{Four levels of measurement:}
			\begin{itemize}
				\item \textbf{Nominal:} categories with no order (stocks, colors)
				\item \textbf{Ordinal:} categories with order, no equal intervals (education, Likert)
				\item \textbf{Interval:} equal intervals, no true zero (temperature, IQ)
				\item \textbf{Ratio:} equal intervals, true zero (income, age, height)
			\end{itemize}
			\textbf{Exam tip:}
			\begin{itemize}
				\item ``Types of'' = nominal
				\item ``Ranked'' = ordinal
				\item ``Temperature/score'' = interval
				\item ``Amount/count'' = ratio
			\end{itemize}
		\end{block}
	\end{frame}
	
	\begin{frame}{Q6 --- Data Scaling}
		\begin{block}{Question}
			Which statement is correct regarding data scaling?
		\end{block}
		\begin{enumerate}[label=\Alph*.]
			\item Standardization creates a distribution bounded by 0 and 1.
			\item Standardization transforms skewed data into standard normal.
			\item Normalization results in data bounded by -1 and 1.
			\item Normalization preserves correlations between data points.
		\end{enumerate}
		\begin{exampleblock}{Answer: D}
		\end{exampleblock}
		\begin{block}{Explanation}
			Standardization = mean 0, SD 1 (not bounded). Normalization = bounded 0--1. Normalization preserves correlations.
		\end{block}
	\end{frame}
	
	\begin{frame}{Q6 Theory --- Standardization vs Normalization}
		\begin{block}{Theory}
			\textbf{Standardization (Z-score):}
			\begin{itemize}
				\item $z = (x - \mu) / \sigma$
				\item Mean = 0, SD = 1
				\item Not bounded by 0 and 1
				\item Does not change distribution shape
			\end{itemize}
			\textbf{Normalization (Min-Max):}
			\begin{itemize}
				\item $x' = (x - \min) / (\max - \min)$
				\item Bounded by 0 and 1
				\item Preserves correlations
				\item Sensitive to outliers
			\end{itemize}
			\textbf{Exam tip:} ``Bounded 0--1'' = normalization. ``Mean 0, SD 1'' = standardization.
		\end{block}
	\end{frame}
	
	\begin{frame}{Q7 --- Stepwise Regression}
		\begin{block}{Question}
			Based on the AIC table, which model should the analyst select?
		\end{block}
		\begin{center}
			\small
			\begin{tabular}{@{}lll@{}}
				\toprule
				Step & Model & AIC \\
				\midrule
				0 & No predictors & 3202.69 \\
				1 & Age & 3193.36 \\
				1 & Age2 & 3203.33 \\
				1 & Age3 & 3204.66 \\
				1 & Age4 & 3203.63 \\
				1 & Age5 & 3202.20 \\
				2 & Age + Age2 & 3125.50 \\
				2 & Age + Age3 & 3123.20 \\
				2 & Age + Age4 & 3126.69 \\
				2 & Age + Age5 & 3132.87 \\
				3 & Age + Age3 + Age2 & 3125.19 \\
				3 & Age + Age3 + Age4 & 3124.77 \\
				3 & Age + Age3 + Age5 & 3124.50 \\
				\bottomrule
			\end{tabular}
		\end{center}
		\begin{enumerate}[label=\Alph*.]
			\item Age + Age3
			\item Age3
			\item Age + Age5
			\item Age + Age3 + Age5
		\end{enumerate}
		\begin{exampleblock}{Answer: A}
		\end{exampleblock}
		\begin{block}{Explanation}
			Forward stepwise picks lowest AIC. Step 2: Age + Age3 (3123.20). Step 3 adds Age5 but AIC rises, so stop.
		\end{block}
	\end{frame}
	
	\begin{frame}{Q7 Theory --- Stepwise Regression and AIC}
		\begin{block}{Theory}
			\textbf{Stepwise regression:}
			\begin{itemize}
				\item Forward: add one at a time
				\item Backward: remove one at a time
				\item Bidirectional: both
			\end{itemize}
			\textbf{AIC (Akaike Information Criterion):}
			\begin{itemize}
				\item Lower AIC = better model
				\item $AIC = 2k - 2\ln(L)$
				\item $k$ = parameters, $L$ = likelihood
			\end{itemize}
			\textbf{Exam tip:} If adding a variable increases AIC, stop.
		\end{block}
	\end{frame}
	
	\begin{frame}{Q8 --- ML vs Classical Statistics}
		\begin{block}{Question}
			Which is a key characteristic of machine learning compared to classical statistics?
		\end{block}
		\begin{enumerate}[label=\Alph*.]
			\item ML focuses on inference and parameter estimation.
			\item ML emphasizes prediction accuracy over interpretability.
			\item ML requires strict assumptions about data distribution.
			\item ML is always more interpretable than classical statistics.
		\end{enumerate}
		\begin{exampleblock}{Answer: B}
		\end{exampleblock}
		\begin{block}{Explanation}
			ML prioritizes prediction. Classical statistics focuses on inference and assumptions.
		\end{block}
	\end{frame}
	
	\begin{frame}{Q8 Theory --- ML vs Classical Statistics}
		\begin{block}{Theory}
			\textbf{Classical statistics:}
			\begin{itemize}
				\item Goal: inference (understand relationships)
				\item Parameter estimation, confidence intervals
				\item Strict assumptions (normality, independence)
				\item More interpretable
			\end{itemize}
			\textbf{Machine learning:}
			\begin{itemize}
				\item Goal: prediction (accuracy on unseen data)
				\item Generalization over inference
				\item Fewer assumptions
				\item Often less interpretable (black-box)
			\end{itemize}
			\textbf{Exam tip:} ``Inference'' = classical stats. ``Prediction'' = ML.
		\end{block}
	\end{frame}
	
	\begin{frame}{Q9 --- Feature Scaling}
		\begin{block}{Question}
			Features include age (20--80) and income (20,000--500,000). Which technique should be applied?
		\end{block}
		\begin{enumerate}[label=\Alph*.]
			\item Remove the income feature
			\item Apply feature scaling
			\item Add more features
			\item Reduce observations
		\end{enumerate}
		\begin{exampleblock}{Answer: B}
		\end{exampleblock}
		\begin{block}{Explanation}
			Feature scaling prevents large-scale features from dominating distance calculations.
		\end{block}
	\end{frame}
	
	\begin{frame}{Q9 Theory --- Feature Scaling}
		\begin{block}{Theory}
			\textbf{Why scale?}
			\begin{itemize}
				\item Distance-based algorithms use magnitudes
				\item Income (20k--500k) vs age (20--80) --- income dominates
			\end{itemize}
			\textbf{When to scale:}
			\begin{itemize}
				\item Distance-based: K-means, KNN, SVM, PCA
				\item Gradient-based: neural networks, logistic regression
				\item Not needed: decision trees, random forests
			\end{itemize}
			\textbf{Methods:} Standardization (mean 0, SD 1) or Normalization (0--1).
		\end{block}
	\end{frame}
	
	\begin{frame}{Q10 --- PCA}
		\begin{block}{Question}
			What is the primary purpose of PCA?
		\end{block}
		\begin{enumerate}[label=\Alph*.]
			\item Increase number of features
			\item Reduce dimensionality while preserving variance
			\item Classify data
			\item Cluster data
		\end{enumerate}
		\begin{exampleblock}{Answer: B}
		\end{exampleblock}
		\begin{block}{Explanation}
			PCA reduces dimensionality by transforming features into uncorrelated components while preserving variance.
		\end{block}
	\end{frame}
	
	\begin{frame}{Q10 Theory --- Principal Component Analysis}
		\begin{block}{Theory}
			\textbf{PCA:}
			\begin{itemize}
				\item Dimensionality reduction technique
				\item Creates uncorrelated components
				\item Ordered by variance explained
				\item First component = most variance
			\end{itemize}
			\textbf{Benefits:}
			\begin{itemize}
				\item Reduces overfitting
				\item Speeds up training
				\item Handles multicollinearity
				\item Enables visualization
			\end{itemize}
			\textbf{Limitations:}
			\begin{itemize}
				\item Hard to interpret components
				\item Assumes linear relationships
				\item Sensitive to scaling
			\end{itemize}
			\textbf{Exam tip:} PCA is unsupervised.
		\end{block}
	\end{frame}
	
	% ============================================================
	\section{Module 2: Tools and Techniques}
	% ============================================================
	
	\begin{frame}{Q11 --- Choosing K}
		\begin{block}{Question}
			How would you choose the number of clusters in unsupervised learning?
		\end{block}
		\begin{enumerate}[label=\Alph*.]
			\item Always use K = 2
			\item Use a scree plot and look for the elbow
			\item Set K equal to number of data points
			\item Use highest possible K
		\end{enumerate}
		\begin{exampleblock}{Answer: B}
		\end{exampleblock}
		\begin{block}{Explanation}
			Scree plot: chart WCSS vs K. Elbow = optimal K. Silhouette scores can also be used.
		\end{block}
	\end{frame}
	
	\begin{frame}{Q11 Theory --- Scree Plot and Silhouette}
		\begin{block}{Theory}
			\textbf{Scree plot (elbow method):}
			\begin{itemize}
				\item Plot WCSS (inertia) vs K
				\item Look for elbow where decrease sharply changes
				\item Beyond elbow, diminishing returns
			\end{itemize}
			\textbf{Silhouette score:}
			\begin{itemize}
				\item $s = (b - a) / \max(a, b)$
				\item $a$ = avg distance to own cluster
				\item $b$ = avg distance to nearest other cluster
				\item Range $-1$ to 1; higher = better
			\end{itemize}
			\textbf{Rule of thumb:} $K \approx \sqrt{N/2}$
		\end{block}
	\end{frame}
	
	\begin{frame}{Q12 --- K-means Steps}
		\begin{block}{Question}
			What are the steps in the K-means clustering algorithm?
		\end{block}
		\begin{enumerate}[label=\Alph*.]
			\item Specify K, scale, select centroids, assign, recalculate, repeat
			\item Specify K, assign labels, train classifier, evaluate
			\item Select K points, remove outliers, apply PCA, cluster
			\item Normalize, hierarchical cluster, cut dendrogram
		\end{enumerate}
		\begin{exampleblock}{Answer: A}
		\end{exampleblock}
		\begin{block}{Explanation}
			K-means: (1) Specify K, (2) Scale features, (3) Random centroids, (4) Assign points, (5) Recalculate centroids, (6) Repeat until stable.
		\end{block}
	\end{frame}
	

	
	\begin{frame}{Q13 --- Initial Centroids}
		\begin{block}{Question}
			How would you choose between clusters from different initial centroid choices?
		\end{block}
		\begin{enumerate}[label=\Alph*.]
			\item Highest WCSS
			\item Lowest WCSS
			\item Most clusters
			\item First solution found
		\end{enumerate}
		\begin{exampleblock}{Answer: B}
		\end{exampleblock}
		\begin{block}{Explanation}
			Lowest WCSS = best fit to feature data.
		\end{block}
	\end{frame}
	
	\begin{frame}{Q13 Theory --- WCSS and BCSS}
		\begin{block}{Theory}
			\textbf{WCSS (Within-Cluster Sum of Squares):}
			\begin{itemize}
				\item Measures compactness
				\item Sum of squared distances to centroid
				\item Lower = better
				\item Always decreases as K increases
			\end{itemize}
			\textbf{BCSS (Between-Cluster Sum of Squares):}
			\begin{itemize}
				\item Measures separation
				\item Sum of squared distances between cluster centroids and overall centroid
				\item Higher = better separation
			\end{itemize}
			\textbf{Exam tip:} WCSS = compactness, BCSS = separation.
		\end{block}
	\end{frame}
	
	\begin{frame}{Q14 --- Scree Plot Use}
		\begin{block}{Question}
			How do you use a scree plot for K-means?
		\end{block}
		\begin{enumerate}[label=\Alph*.]
			\item Highest point
			\item Lowest point
			\item Elbow where gradient changes from steep to flat
			\item Point where WCSS equals zero
		\end{enumerate}
		\begin{exampleblock}{Answer: C}
		\end{exampleblock}
		\begin{block}{Explanation}
			The elbow = optimal K. Beyond it, adding clusters gives diminishing returns.
		\end{block}
	\end{frame}
	
	\begin{frame}{Q14 Theory --- Scree Plot Interpretation}
		\begin{block}{Theory}
			\textbf{Scree plot:}
			\begin{itemize}
				\item X-axis: K
				\item Y-axis: WCSS
				\item Always decreases
				\item Look for elbow (max curvature)
			\end{itemize}
			\textbf{Elbow:}
			\begin{itemize}
				\item Before: steep decrease
				\item After: flat decrease
				\item Optimal trade-off
			\end{itemize}
			\textbf{Pitfalls:}
			\begin{itemize}
				\item Ambiguous elbow
				\item No visible elbow
				\item Use domain knowledge + silhouette
			\end{itemize}
		\end{block}
	\end{frame}
	
	\begin{frame}{Q15 --- Hierarchical Clustering}
		\begin{block}{Question}
			What are the two types of hierarchical clustering and their advantages?
		\end{block}
		\begin{enumerate}[label=\Alph*.]
			\item Divisive and agglomerative; no need to pre-define K, dendrograms interpretable
			\item Supervised and unsupervised; faster, more accurate
			\item Partitional and density-based; handles outliers
			\item Single and complete linkage; simpler, more scalable
		\end{enumerate}
		\begin{exampleblock}{Answer: A}
		\end{exampleblock}
		\begin{block}{Explanation}
			Divisive (top-down) and agglomerative (bottom-up). No need to pre-define K. Dendrograms are interpretable.
		\end{block}
	\end{frame}
	
	\begin{frame}{Q15 Theory --- Hierarchical Clustering}
		\begin{block}{Theory}
			\textbf{Approaches:}
			\begin{itemize}
				\item Divisive: top-down splitting
				\item Agglomerative: bottom-up merging
			\end{itemize}
			\textbf{Linkage criteria:}
			\begin{itemize}
				\item Single: shortest distance
				\item Complete: greatest distance
				\item Average: average distance
				\item Ward's: minimizes variance
			\end{itemize}
			\textbf{Advantages:}
			\begin{itemize}
				\item No need to pre-define K
				\item Dendrogram shows all clusterings
				\item Uncovers nested relationships
			\end{itemize}
		\end{block}
	\end{frame}
	
	\begin{frame}{Q16 --- K-means Spherical Clusters}
		\begin{block}{Question}
			True or False: K-means with Euclidean distance can only be used when clusters are approximately spherical.
		\end{block}
		\begin{enumerate}[label=\Alph*.]
			\item True
			\item False
			\item True only for large datasets
			\item True only for small datasets
		\end{enumerate}
		\begin{exampleblock}{Answer: A}
		\end{exampleblock}
		\begin{block}{Explanation}
			True. K-means assumes spherical clusters. Manhattan distance produces rhombus-shaped clusters.
		\end{block}
	\end{frame}
	
	\begin{frame}{Q16 Theory --- K-means Limitations}
		\begin{block}{Theory}
			\textbf{K-means produces spherical clusters:}
			\begin{itemize}
				\item Euclidean distance $\rightarrow$ spherical
				\item Manhattan distance $\rightarrow$ rhombus
			\end{itemize}
			\textbf{Problems:}
			\begin{itemize}
				\item Outliers distort centroids
				\item Curse of dimensionality
				\item Non-spherical clusters fail
				\item Correlated features struggle
			\end{itemize}
			\textbf{Solutions:}
			\begin{itemize}
				\item Standardize, remove outliers
				\item Apply PCA
				\item Use kernel functions or DBSCAN
			\end{itemize}
		\end{block}
	\end{frame}
	
	\begin{frame}{Q17 --- WCSS Monotonicity}
		\begin{block}{Question}
			True or False: WCSS will never rise when K increases.
		\end{block}
		\begin{enumerate}[label=\Alph*.]
			\item True
			\item False
			\item True only if K is even
			\item True only if K is odd
		\end{enumerate}
		\begin{exampleblock}{Answer: A}
		\end{exampleblock}
		\begin{block}{Explanation}
			True. Adding clusters cannot increase WCSS. At K = N, WCSS = 0.
		\end{block}
	\end{frame}
	
	\begin{frame}{Q17 Theory --- WCSS Monotonicity}
		\begin{block}{Theory}
			\textbf{WCSS is monotonically non-increasing in K:}
			\begin{itemize}
				\item Adding cluster never increases WCSS
				\item At K = N, WCSS = 0
				\item But K = N is overfitting
			\end{itemize}
			\textbf{Analogy:}
			\begin{itemize}
				\item Regression: adding features never decreases $R^2$
				\item K-means: adding clusters never increases WCSS
			\end{itemize}
			\textbf{Implication:}
			\begin{itemize}
				\item Cannot choose K by WCSS alone
				\item Use elbow, silhouette, or domain knowledge
			\end{itemize}
		\end{block}
	\end{frame}
	
	\begin{frame}{Q18 --- Dendrogram}
		\begin{block}{Question}
			What is a dendrogram and how is it interpreted?
		\end{block}
		\begin{enumerate}[label=\Alph*.]
			\item Bar chart of feature importance
			\item Tree-like diagram; long vertical lines = important clusters
			\item Scatter plot of clusters
			\item Line chart of WCSS vs K
		\end{enumerate}
		\begin{exampleblock}{Answer: B}
		\end{exampleblock}
		\begin{block}{Explanation}
			Dendrogram = tree showing hierarchical clustering. Long vertical lines = important cluster boundaries.
		\end{block}
	\end{frame}
	
	\begin{frame}{Q18 Theory --- Dendrograms}
		\begin{block}{Theory}
			\textbf{Dendrogram:}
			\begin{itemize}
				\item X-axis: data points
				\item Y-axis: distance
				\item Each merge = horizontal line
				\item Height = merge distance
			\end{itemize}
			\textbf{Interpretation:}
			\begin{itemize}
				\item Long vertical lines = important clusters
				\item Short lines = minor clusters
				\item Cut horizontally to choose K
				\item Number of vertical lines crossed = K
			\end{itemize}
		\end{block}
	\end{frame}
	
	\begin{frame}{Q19 --- Centroid Calculation}
		\begin{block}{Question}
			Data on three features for banks A, B, C. What is the centroid?
			\begin{center}
				\small
				\begin{tabular}{@{}lccc@{}}
					\toprule
					Feature & A & B & C \\
					\midrule
					Customers (M) & 1.2 & 6.0 & 0.5 \\
					Loan book (bn) & 5 & 25 & 7 \\
					Branches & 80 & 400 & 50 \\
					\bottomrule
				\end{tabular}
			\end{center}
		\end{block}
		\begin{enumerate}[label=\Alph*.]
			\item (2.567, 12.333, 176.667)
			\item (2.567, 12.333, 176.667)
			\item (2.567, 12.333, 176.667)
			\item (2.567, 12.333, 176.667)
		\end{enumerate}
		\begin{exampleblock}{Answer: A}
		\end{exampleblock}
		\begin{block}{Explanation}
			Average: (1.2+6.0+0.5)/3 = 2.567; (5+25+7)/3 = 12.333; (80+400+50)/3 = 176.667.
		\end{block}
	\end{frame}
	
	\begin{frame}{Q19 Theory --- Centroids}
		\begin{block}{Theory}
			\textbf{Centroid = mean of cluster points:}
			\begin{itemize}
				\item Average each feature across all points
				\item Recalculated each K-means iteration
				\item Does not have to be an actual data point
			\end{itemize}
			\textbf{Formula:}
			\[
			\mu_j = \frac{1}{n_k} \sum_{i \in C_k} x_i
			\]
			\textbf{Exam tip:} Add and divide by count.
		\end{block}
	\end{frame}
	
	\begin{frame}{Q20 --- K-means Issues}
		\begin{block}{Question}
			Which issue might pose problems for K-means?
		\end{block}
		\begin{enumerate}[label=\Alph*.]
			\item Outliers
			\item High dimensionality
			\item Non-linearity
			\item Non-spherical clusters
		\end{enumerate}
		\begin{exampleblock}{Answer: D}
		\end{exampleblock}
		\begin{block}{Explanation}
			K-means fails with non-spherical clusters. It assumes spherical clusters.
		\end{block}
	\end{frame}
	
	\begin{frame}{Q20 Theory --- K-means Failure Modes}
		\begin{block}{Theory}
			\textbf{Two pathological cases:}
			\begin{itemize}
				\item Concentric circles: centroids end up in same place
				\item Elongated clusters: edge points misallocated
			\end{itemize}
			\textbf{Why K-means fails:}
			\begin{itemize}
				\item Assumes spherical, similar-size clusters
				\item Uses distance to centroid only
			\end{itemize}
			\textbf{Solutions:}
			\begin{itemize}
				\item Kernel K-means
				\item DBSCAN
				\item Spectral clustering
				\item Gaussian Mixture Models
			\end{itemize}
		\end{block}
	\end{frame}
	
	\begin{frame}{Q21 --- DBSCAN}
		\begin{block}{Question}
			What is the main advantage of DBSCAN over K-means?
		\end{block}
		\begin{enumerate}[label=\Alph*.]
			\item DBSCAN is faster
			\item DBSCAN finds arbitrary-shaped clusters and handles outliers
			\item DBSCAN requires fewer parameters
			\item DBSCAN always produces spherical clusters
		\end{enumerate}
		\begin{exampleblock}{Answer: B}
		\end{exampleblock}
		\begin{block}{Explanation}
			DBSCAN finds arbitrary-shaped clusters and is robust to outliers (noise points).
		\end{block}
	\end{frame}
	
	\begin{frame}{Q21 Theory --- DBSCAN}
		\begin{block}{Theory}
			\textbf{DBSCAN:}
			\begin{itemize}
				\item Clusters = dense regions
				\item Core point: at least minPts neighbors within radius $\varepsilon$
				\item Border point: within $\varepsilon$ of core, not core itself
				\item Noise point: neither = outlier
			\end{itemize}
			\textbf{Advantages:}
			\begin{itemize}
				\item Arbitrary shapes
				\item Robust to outliers
				\item No need for K
			\end{itemize}
			\textbf{Disadvantages:}
			\begin{itemize}
				\item Varying densities
				\item Sensitive to $\varepsilon$ and minPts
			\end{itemize}
		\end{block}
	\end{frame}
	
	\begin{frame}{Q22 --- Silhouette Score}
		\begin{block}{Question}
			What is the purpose of the silhouette score?
		\end{block}
		\begin{enumerate}[label=\Alph*.]
			\item Measure accuracy
			\item Compare cluster cohesion and separation
			\item Count clusters
			\item Visualize high-dimensional data
		\end{enumerate}
		\begin{exampleblock}{Answer: B}
		\end{exampleblock}
		\begin{block}{Explanation}
			Compares distance to own cluster (cohesion) with distance to nearest other cluster (separation).
		\end{block}
	\end{frame}
	
	\begin{frame}{Q22 Theory --- Silhouette Score}
		\begin{block}{Theory}
			\textbf{Silhouette score for point $i$:}
			\[
			s_i = \frac{b_i - a_i}{\max(b_i, a_i)}
			\]
			\begin{itemize}
				\item $a_i$ = mean distance to own cluster
				\item $b_i$ = mean distance to nearest other cluster
				\item Range: $-1$ to 1
			\end{itemize}
			\textbf{Interpretation:}
			\begin{itemize}
				\item $s \approx 1$: well inside cluster
				\item $s \approx 0$: boundary
				\item $s < 0$: possibly mis-clustered
			\end{itemize}
			\textbf{Choose K with highest average silhouette.}
		\end{block}
	\end{frame}
	
	\begin{frame}{Q23 --- Curse of Dimensionality}
		\begin{block}{Question}
			What is the curse of dimensionality in K-means?
		\end{block}
		\begin{enumerate}[label=\Alph*.]
			\item Algorithm too slow with many observations
			\item Data points become sparse in high dimensions
			\item Number of clusters increases exponentially
			\item Centroids become unstable
		\end{enumerate}
		\begin{exampleblock}{Answer: B}
		\end{exampleblock}
		\begin{block}{Explanation}
			As features increase, data becomes sparse. Distances become similar, making clustering hard.
		\end{block}
	\end{frame}
	
	\begin{frame}{Q23 Theory --- Curse of Dimensionality}
		\begin{block}{Theory}
			\textbf{Curse of dimensionality:}
			\begin{itemize}
				\item Distances become similar with more features
				\item K-means cannot distinguish clusters
				\item Convergence slows
			\end{itemize}
			\textbf{Solutions:}
			\begin{itemize}
				\item Remove irrelevant features
				\item Apply PCA
				\item Use cosine distance
				\item Use algorithms designed for high dimensions
			\end{itemize}
		\end{block}
	\end{frame}
	
	\begin{frame}{Q24 --- Linkage}
		\begin{block}{Question}
			What is the difference between single linkage and complete linkage?
		\end{block}
		\begin{enumerate}[label=\Alph*.]
			\item Single uses shortest distance; complete uses greatest
			\item Single uses greatest; complete uses shortest
			\item Single for agglomerative; complete for divisive
			\item Single for divisive; complete for agglomerative
		\end{enumerate}
		\begin{exampleblock}{Answer: A}
		\end{exampleblock}
		\begin{block}{Explanation}
			Single linkage = shortest distance between clusters. Complete = greatest distance.
		\end{block}
	\end{frame}
	
	\begin{frame}{Q24 Theory --- Linkage Criteria}
		\begin{block}{Theory}
			\textbf{Linkage = cluster distance measure:}
			\begin{itemize}
				\item Single: shortest distance between any two points
				\item Complete: greatest distance
				\item Average: average distance
				\item Ward's: minimizes within-cluster variance
			\end{itemize}
			\textbf{Properties:}
			\begin{itemize}
				\item Single: elongated clusters, sensitive to outliers
				\item Complete: compact clusters, robust to outliers
				\item Average: balance
			\end{itemize}
			\textbf{Exam tip:} Single = shortest. Complete = greatest.
		\end{block}
	\end{frame}
	
	\begin{frame}{Q25 --- K-means++}
		\begin{block}{Question}
			What is the purpose of K-means++ initialization?
		\end{block}
		\begin{enumerate}[label=\Alph*.]
			\item Speed up convergence by choosing centroids far apart
			\item Reduce number of clusters
			\item Increase WCSS
			\item Eliminate outliers
		\end{enumerate}
		\begin{exampleblock}{Answer: A}
		\end{exampleblock}
		\begin{block}{Explanation}
			K-means++ spreads initial centroids apart, reducing poor clustering and speeding convergence.
		\end{block}
	\end{frame}
	
	
	
	
	\begin{frame}{Q26 --- Convolution}
		\begin{block}{Question}
			Given X and W, what is the feature map F?
			\[
			X = \begin{pmatrix} 0 & 1 & 2 \\ 1 & 0 & 0 \\ 3 & 2 & 2 \end{pmatrix}, \quad
			W = \begin{pmatrix} -1 & 0 \\ 1 & 1 \end{pmatrix}
			\]
		\end{block}
		\begin{enumerate}[label=\Alph*.]
			\item $F = \begin{pmatrix} 1 & -1 \\ 4 & 4 \end{pmatrix}$
			\item $F = \begin{pmatrix} 1 & 4 \\ -1 & 4 \end{pmatrix}$
			\item $F = \begin{pmatrix} 2 & 3 \\ 6 & 4 \end{pmatrix}$
			\item $F = \begin{pmatrix} 2 & 6 \\ 3 & 4 \end{pmatrix}$
		\end{enumerate}
		\begin{exampleblock}{Answer: A}
		\end{exampleblock}
		\begin{block}{Explanation}
			F11 = 0(-1)+1(0)+1(1)+0(1) = 1; F12 = -1; F21 = 4; F22 = 4.
		\end{block}
	\end{frame}
	
	\begin{frame}{Q26 Theory --- Convolution}
		\begin{block}{Theory}
			\textbf{Convolution:}
			\begin{itemize}
				\item Slide kernel over input
				\item Multiply overlapping, sum
				\item Result = feature map
			\end{itemize}
			\textbf{Example (2x2 kernel on 3x3):}
			\begin{itemize}
				\item Output size = $(3-2+1) \times (3-2+1) = 2 \times 2$
				\item $F_{11} = X_{11}W_{11} + X_{12}W_{12} + X_{21}W_{21} + X_{22}W_{22}$
			\end{itemize}
			\textbf{Used in:} CNNs for image recognition.
		\end{block}
	\end{frame}
	
	\begin{frame}{Q27 --- Ensemble}
		\begin{block}{Question}
			Analyst samples with replacement, creates multiple trees using all features, averages. Which technique?
		\end{block}
		\begin{enumerate}[label=\Alph*.]
			\item Adaptive boosting
			\item Random forest
			\item Gradient boosting
			\item Bootstrap aggregation
		\end{enumerate}
		\begin{exampleblock}{Answer: D}
		\end{exampleblock}
		\begin{block}{Explanation}
			Bagging uses all features. Random forest uses random subset of features.
		\end{block}
	\end{frame}
	
	\begin{frame}{Q27 Theory --- Ensemble Techniques}
		\begin{block}{Theory}
			\textbf{Ensemble = combine models:}
			\begin{itemize}
				\item \textbf{Bagging:} bootstrap samples, all features, average
				\item \textbf{Random Forest:} bagging + random feature subset
				\item \textbf{Boosting:} sequential, correct errors
				\begin{itemize}
					\item AdaBoost: reweights misclassified
					\item Gradient Boosting: fits residuals
				\end{itemize}
			\end{itemize}
			\textbf{Exam tip:}
			\begin{itemize}
				\item ``All features'' + ``average'' = bagging
				\item ``Random subset'' = random forest
				\item ``Sequential'' = boosting
			\end{itemize}
		\end{block}
	\end{frame}
	
	\begin{frame}{Q28 --- Random Forest Advantage}
		\begin{block}{Question}
			What is an advantage of random forest?
		\end{block}
		\begin{enumerate}[label=\Alph*.]
			\item Captures previous misclassifications
			\item Improves interpretability
			\item Improves accuracy by reducing correlation across trees
			\item Captures outliers via Cook's distance
		\end{enumerate}
		\begin{exampleblock}{Answer: C}
		\end{exampleblock}
		\begin{block}{Explanation}
			Random forest reduces correlation across trees, improving prediction accuracy.
		\end{block}
	\end{frame}
	
	\begin{frame}{Q28 Theory --- Random Forest}
		\begin{block}{Theory}
			\textbf{Random Forest:}
			\begin{itemize}
				\item Each tree on bootstrap sample\item Random feature subset at each split
				\item Average (regression) or majority vote (classification)
			\end{itemize}
			\textbf{Advantages:}
			\begin{itemize}
				\item Reduces overfitting
				\item Reduces correlation across trees
				\item Handles high dimensions
				\item Robust to outliers
				\item Feature importance
			\end{itemize}
			\textbf{Disadvantages:} Less interpretable, expensive.
		\end{block}
	\end{frame}
	
	\begin{frame}{Q29 --- SVM Classification}
		\begin{block}{Question}
			Boundary: $0.43x_1 + 2.55x_2 + 0.67 = 0$. For $x_1 = 0.8$, $x_2 = -0.75$, what class?
		\end{block}
		\begin{enumerate}[label=\Alph*.]
			\item Class 0, decision value lower than 0
			\item Class 0, decision value higher than 0
			\item Class 1, decision value lower than 0
			\item Class 1, decision value higher than 0
		\end{enumerate}
		\begin{exampleblock}{Answer: A}
		\end{exampleblock}
		\begin{block}{Explanation}
			$0.43(0.8) + 2.55(-0.75) + 0.67 = -0.8985 < 0 \rightarrow$ Class 0.
		\end{block}
	\end{frame}
	
	\begin{frame}{Q29 Theory --- SVM}
		\begin{block}{Theory}
			\textbf{SVM:}
			\begin{itemize}
				\item Boundary: $w_1x_1 + w_2x_2 + b = 0$
				\item $f(x) > 0 \rightarrow$ Class 1
				\item $f(x) < 0 \rightarrow$ Class 0
			\end{itemize}
			\textbf{Key concepts:}
			\begin{itemize}
				\item Support vectors: closest points
				\item Margin: distance to support vectors
				\item Kernel trick: higher dimensions
			\end{itemize}
			\textbf{Exam tip:} Positive = Class 1, negative = Class 0.
		\end{block}
	\end{frame}
	
	\begin{frame}{Q30 --- Classification Metrics}
		\begin{block}{Question}
			Which model correctly identifies the highest percent of defaulted loans?
			\begin{center}
				\small
				\begin{tabular}{@{}lcccc@{}}
					\toprule
					Model & Acc & Prec & Rec & F1 \\
					\midrule
					A (Logistic) & 0.71 & 0.60 & 0.28 & 0.38 \\
					B (Tree) & 0.66 & 0.43 & 0.17 & 0.24 \\
					C (SVM) & 0.63 & 0.44 & 0.51 & 0.47 \\
					D (KNN) & 0.68 & 0.50 & 0.15 & 0.23 \\
					\bottomrule
				\end{tabular}
			\end{center}
		\end{block}
		\begin{enumerate}[label=\Alph*.]
			\item Model A
			\item Model B
			\item Model C
			\item Model D
		\end{enumerate}
		\begin{exampleblock}{Answer: C}
		\end{exampleblock}
		\begin{block}{Explanation}
			Model C has highest recall (0.51) = true positive rate.
		\end{block}
	\end{frame}
	
	\begin{frame}{Q30 Theory --- Classification Metrics}
		\begin{block}{Theory}
			\textbf{Confusion matrix:}
			\begin{center}
				\small
				\begin{tabular}{@{}lcc@{}}
					\toprule
					& Predicted + & Predicted - \\
					\midrule
					Actual + & TP & FN \\
					Actual - & FP & TN \\
					\bottomrule
				\end{tabular}
			\end{center}
			\textbf{Metrics:}
			\begin{itemize}
				\item Accuracy = (TP+TN)/Total
				\item Precision = TP/(TP+FP)
				\item Recall = TP/(TP+FN)
				\item F1 = 2 x Prec x Rec / (Prec + Rec)
				\item Specificity = TN/(TN+FP)
			\end{itemize}
			\textbf{Exam tip:} ``Percent of defaulted loans caught'' = recall.
		\end{block}
	\end{frame}
	
	\begin{frame}{Q31 --- Tokenization}
		\begin{block}{Question}
			What does tokenization do?
		\end{block}
		\begin{enumerate}[label=\Alph*.]
			\item Removes words with no value
			\item Separates text into words
			\item Converts words to numeric vectors
			\item Replaces words with similar meanings
		\end{enumerate}
		\begin{exampleblock}{Answer: B}
		\end{exampleblock}
		\begin{block}{Explanation}
			Tokenization = split text into words/tokens.
		\end{block}
	\end{frame}
	

	
	\begin{frame}{Q32 --- CBoW}
		\begin{block}{Question}
			What does CBoW represent?
		\end{block}
		\begin{enumerate}[label=\Alph*.]
			\item Recurrent neural network
			\item Fill-in-the-blank technique predicting a missing word
			\item Skip-gram
			\item Auto-encoder
		\end{enumerate}
		\begin{exampleblock}{Answer: B}
		\end{exampleblock}
		\begin{block}{Explanation}
			CBoW = fill-in-the-blank. Predicts missing word from context.
		\end{block}
	\end{frame}
	
	\begin{frame}{Q32 Theory --- Word2Vec Architectures}
		\begin{block}{Theory}
			\textbf{Word2Vec:}
			\begin{itemize}
				\item Converts words to dense vectors
				\item Captures semantic relationships
			\end{itemize}
			\textbf{CBoW:}
			\begin{itemize}
				\item Predicts missing word from context
				\item Multiple inputs, one output
			\end{itemize}
			\textbf{Skip-gram:}
			\begin{itemize}
				\item Predicts context from target word
				\item One input, multiple outputs
				\item Better with rare words
			\end{itemize}
			\textbf{Exam tip:} CBoW = many in, one out. Skip-gram = one in, many out.
		\end{block}
	\end{frame}
	
	\begin{frame}{Q33 --- Co-training}
		\begin{block}{Question}
			What is an advantage of co-training?
		\end{block}
		\begin{enumerate}[label=\Alph*.]
			\item Different views of data; models learn from each other
			\item Higher productivity via parallel computing
			\item Improves accuracy by decreasing weights
			\item More reliable; data not interchangeable
		\end{enumerate}
		\begin{exampleblock}{Answer: A}
		\end{exampleblock}
		\begin{block}{Explanation}
			Co-training provides different views; models learn from each other.
		\end{block}
	\end{frame}
	
	\begin{frame}{Q33 Theory --- Semi-Supervised Learning}
		\begin{block}{Theory}
			\textbf{Semi-supervised:} small labeled + large unlabeled.
			\vspace{2mm}
			
			\textbf{Techniques:}
			\begin{itemize}
				\item Self-training: predict unlabeled, add confident, repeat
				\item Co-training: split features, two models label for each other
				\item Unsupervised pre-processing: cluster then classify
			\end{itemize}
			\textbf{Assumptions:}
			\begin{itemize}
				\item Smoothness: nearby points similar
				\item Cluster: same cluster same label
				\item Manifold: lower-dimensional
			\end{itemize}
		\end{block}
	\end{frame}
	
	\begin{frame}{Q34 --- Regularization}
		\begin{block}{Question}
			What is the primary purpose of regularization?
		\end{block}
		\begin{enumerate}[label=\Alph*.]
			\item Increase model complexity
			\item Prevent overfitting by penalizing large coefficients
			\item Speed up training
			\item Increase features
		\end{enumerate}
		\begin{exampleblock}{Answer: B}
		\end{exampleblock}
		\begin{block}{Explanation}
			Regularization (Ridge, LASSO, Elastic Net) prevents overfitting.
		\end{block}
	\end{frame}
	
	\begin{frame}{Q34 Theory --- Regularization}
		\begin{block}{Theory}
			\textbf{Regularization:}
			\begin{itemize}
				\item Ridge (L2): $\lambda \sum \beta_j^2$; shrinks, keeps all
				\item LASSO (L1): $\lambda \sum |\beta_j|$; can zero coefficients
				\item Elastic Net: both
			\end{itemize}
			\textbf{Exam tip:}
			\begin{itemize}
				\item Feature selection = LASSO
				\item Shrink coefficients = Ridge
				\item Both = Elastic Net
			\end{itemize}
		\end{block}
	\end{frame}
	
	\begin{frame}{Q35 --- Ridge vs LASSO}
		\begin{block}{Question}
			What is the difference between Ridge and LASSO?
		\end{block}
		\begin{enumerate}[label=\Alph*.]
			\item Ridge L1; LASSO L2
			\item Ridge L2; LASSO L1
			\item Both L1
			\item Both L2
		\end{enumerate}
		\begin{exampleblock}{Answer: B}
		\end{exampleblock}
		\begin{block}{Explanation}
			Ridge = L2 (squared). LASSO = L1 (absolute), can zero coefficients.
		\end{block}
	\end{frame}
	
	\begin{frame}{Q35 Theory --- Ridge vs LASSO}
		\begin{block}{Theory}
			\textbf{Ridge (L2):}
			\[
			\text{Loss} = \text{RSS} + \lambda \sum_{j=1}^{p} \beta_j^2
			\]
			Shrinks toward zero, never exactly zero. Keeps all features.
			\vspace{2mm}
			
			\textbf{LASSO (L1):}
			\[
			\text{Loss} = \text{RSS} + \lambda \sum_{j=1}^{p} |\beta_j|
			\]
			Can zero coefficients. Feature selection.
			\vspace{2mm}
			
			\textbf{Elastic Net:} both penalties.
		\end{block}
	\end{frame}
	
	\begin{frame}{Q36 --- Cross-Validation}
		\begin{block}{Question}
			What is the purpose of cross-validation?
		\end{block}
		\begin{enumerate}[label=\Alph*.]
			\item Increase training size
			\item Estimate model performance on unseen data
			\item Reduce features
			\item Speed up training
		\end{enumerate}
		\begin{exampleblock}{Answer: B}
		\end{exampleblock}
		\begin{block}{Explanation}
			Cross-validation estimates generalization to unseen data.
		\end{block}
	\end{frame}
	
	\begin{frame}{Q36 Theory --- Cross-Validation}
		\begin{block}{Theory}
			\textbf{Cross-validation:}
			\begin{itemize}
				\item K-fold: train on K-1, validate on 1, repeat K times
				\item Stratified: preserve class proportions
				\item LOO: K = N
				\item Time-series: respect temporal order
			\end{itemize}
			\textbf{Benefits:}
			\begin{itemize}
				\item Reliable estimate
				\item Uses all data
				\item Detects overfitting
			\end{itemize}
			\textbf{Exam tip:} CV is for evaluation, not training.
		\end{block}
	\end{frame}
	
	\begin{frame}{Q37 --- Bias-Variance}
		\begin{block}{Question}
			What is the bias-variance tradeoff?
		\end{block}
		\begin{enumerate}[label=\Alph*.]
			\item Accuracy vs training time
			\item Underfitting vs overfitting
			\item Features vs observations
			\item Precision vs recall
		\end{enumerate}
		\begin{exampleblock}{Answer: B}
		\end{exampleblock}
		\begin{block}{Explanation}
			High bias = underfitting. High variance = overfitting.
		\end{block}
	\end{frame}
	
	\begin{frame}{Q37 Theory --- Bias-Variance Tradeoff}
		\begin{block}{Theory}
			\textbf{Total error = Bias$^2$ + Variance + Irreducible Error.}
			\vspace{2mm}
			
			\textbf{Bias:} overly simple assumptions (underfitting).\\
			\textbf{Variance:} sensitivity to training data (overfitting).
			\vspace{2mm}
			
			\textbf{Tradeoff:}
			\begin{itemize}
				\item Simple: high bias, low variance
				\item Complex: low bias, high variance
				\item Optimal: balance
			\end{itemize}
			\textbf{Diagnose:}
			\begin{itemize}
				\item High train + high test error = underfitting
				\item Low train + high test error = overfitting
			\end{itemize}
		\end{block}
	\end{frame}
	
	\begin{frame}{Q38 --- Confusion Matrix}
		\begin{block}{Question}
			What is the purpose of a confusion matrix?
		\end{block}
		\begin{enumerate}[label=\Alph*.]
			\item Visualize decision boundary
			\item Summarize classification performance
			\item Plot ROC curve
			\item Cluster data
		\end{enumerate}
		\begin{exampleblock}{Answer: B}
		\end{exampleblock}
		\begin{block}{Explanation}
			Confusion matrix shows TP, TN, FP, FN.
		\end{block}
	\end{frame}
	
	\begin{frame}{Q38 Theory --- Confusion Matrix}
		\begin{block}{Theory}
			\textbf{Confusion matrix:}
			\begin{center}
				\small
				\begin{tabular}{@{}lcc@{}}
					\toprule
					& Predicted + & Predicted - \\
					\midrule
					Actual + & TP & FN \\
					Actual - & FP & TN \\
					\bottomrule
				\end{tabular}
			\end{center}
			\textbf{Metrics:}
			\begin{itemize}
				\item Accuracy = (TP+TN)/(TP+TN+FP+FN)
				\item Precision = TP/(TP+FP)
				\item Recall = TP/(TP+FN)
				\item Specificity = TN/(TN+FP)
				\item F1 = 2 x Prec x Rec / (Prec + Rec)
			\end{itemize}
		\end{block}
	\end{frame}
	
	\begin{frame}{Q39 --- AUC-ROC}
		\begin{block}{Question}
			What does the AUC-ROC curve measure?
		\end{block}
		\begin{enumerate}[label=\Alph*.]
			\item Regression accuracy
			\item Classifier's ability to distinguish classes across thresholds
			\item Clustering speed
			\item Number of clusters
		\end{enumerate}
		\begin{exampleblock}{Answer: B}
		\end{exampleblock}
		\begin{block}{Explanation}
			AUC-ROC measures discrimination across all thresholds.
		\end{block}
	\end{frame}
	
	\begin{frame}{Q39 Theory --- ROC and AUC}
		\begin{block}{Theory}
			\textbf{ROC curve:}
			\begin{itemize}
				\item X-axis: FPR = 1 - Specificity
				\item Y-axis: TPR = Sensitivity
				\item Each point = different threshold
			\end{itemize}
			\textbf{AUC:}
			\begin{itemize}
				\item Range 0 to 1
				\item 0.5 = random
				\item 1.0 = perfect
				\item Higher = better
			\end{itemize}
			\textbf{Exam tip:} AUC measures discrimination, not accuracy.
		\end{block}
	\end{frame}
	
	\begin{frame}{Q40 --- Precision vs Recall}
		\begin{block}{Question}
			What is the difference between precision and recall?
		\end{block}
		\begin{enumerate}[label=\Alph*.]
			\item Precision = TP/(TP+FP); Recall = TP/(TP+FN)
			\item Precision = TP/(TP+FN); Recall = TP/(TP+FP)
			\item Same thing
			\item Precision for regression; recall for classification
		\end{enumerate}
		\begin{exampleblock}{Answer: A}
		\end{exampleblock}
		\begin{block}{Explanation}
			Precision = accuracy of positive predictions. Recall = capture of actual positives.
		\end{block}
	\end{frame}
	
	\begin{frame}{Q40 Theory --- Precision vs Recall}
		\begin{block}{Theory}
			\textbf{Precision:} Of predicted positives, how many are correct?
			\[
			\text{Precision} = \frac{TP}{TP + FP}
			\]
			\textbf{Recall:} Of actual positives, how many caught?
			\[
			\text{Recall} = \frac{TP}{TP + FN}
			\]
			\textbf{When:}
			\begin{itemize}
				\item High precision: false positives costly (spam)
				\item High recall: false negatives costly (cancer, fraud)
				\item F1: harmonic mean
			\end{itemize}
		\end{block}
	\end{frame}
	
	% ============================================================
	\section{Module 3: Risk and Risk Factors}
	% ============================================================
	
	\begin{frame}{Q41 --- Rare Disease}
		\begin{block}{Question}
			Rare disease (0.1\%). Untreated = fatal. Which metric?
		\end{block}
		\begin{enumerate}[label=\Alph*.]
			\item Sensitivity
			\item Precision
			\item Accuracy
			\item Specificity
		\end{enumerate}
		\begin{exampleblock}{Answer: A}
		\end{exampleblock}
		\begin{block}{Explanation}
			Sensitivity (true positive rate) because missing cases is fatal.
		\end{block}
	\end{frame}
	
	\begin{frame}{Q41 Theory --- Sensitivity vs Specificity}
		\begin{block}{Theory}
			\textbf{Sensitivity:}
			\[
			\text{Sensitivity} = \frac{TP}{TP + FN}
			\]
			Use when false negatives costly (miss disease).
			\vspace{2mm}
			
			\textbf{Specificity:}
			\[
			\text{Specificity} = \frac{TN}{TN + FP}
			\]
			Use when false positives costly (unnecessary treatment).
			\vspace{2mm}
			
			\textbf{Exam tip:} Rare + fatal = sensitivity.
		\end{block}
	\end{frame}
	
	\begin{frame}{Q42 --- Data Composition Bias}
		\begin{block}{Question}
			10,000 transactions, 20 fraudulent. What bias?
		\end{block}
		\begin{enumerate}[label=\Alph*.]
			\item Measurement bias
			\item Representation bias
			\item Data composition bias
			\item Historical bias
		\end{enumerate}
		\begin{exampleblock}{Answer: C}
		\end{exampleblock}
		\begin{block}{Explanation}
			Imbalanced dataset = data composition bias.
		\end{block}
	\end{frame}
	
	\begin{frame}{Q42 Theory --- Sources of Algorithmic Bias}
		\begin{block}{Theory}
			\textbf{Types:}
			\begin{itemize}
				\item Measurement: features measured differently
				\item Representation: sample not representative
				\item Data composition: imbalanced dataset
				\item Historical: past discrimination encoded
				\item Problem specification: wrong target
				\item Deployment: wrong context
			\end{itemize}
			\textbf{Exam tip:} Imbalanced = data composition bias.
		\end{block}
	\end{frame}
	
	\begin{frame}{Q43 --- Education Level Fairness}
		\begin{block}{Question}
			Best approach to ensure education level does not affect predictions?
		\end{block}
		\begin{enumerate}[label=\Alph*.]
			\item Remove education level
			\item Replace with average during prediction
			\item Add interaction terms
			\item Create dummy variables
		\end{enumerate}
		\begin{exampleblock}{Answer: B}
		\end{exampleblock}
		\begin{block}{Explanation}
			Replacing with average ensures equal treatment.
		\end{block}
	\end{frame}
	
	\begin{frame}{Q43 Theory --- Fairness and Bias Mitigation}
		\begin{block}{Theory}
			\textbf{Approaches:}
			\begin{itemize}
				\item Pre-processing: fix data
				\item In-processing: modify algorithm
				\item Post-processing: adjust predictions
			\end{itemize}
			\textbf{Techniques:}
			\begin{itemize}
				\item Remove sensitive features (proxy effects remain)
				\item Replace with average during prediction
				\item Fairness constraints
				\item Reweight samples
			\end{itemize}
		\end{block}
	\end{frame}
	
	\begin{frame}{Q44 --- Demographic Parity}
		\begin{block}{Question}
			How does demographic parity contribute to group fairness?
		\end{block}
		\begin{enumerate}[label=\Alph*.]
			\item Equal accuracy
			\item Identical predicted positive rates across subgroups
			\item Identical sensitivity
			\item Equal precision
		\end{enumerate}
		\begin{exampleblock}{Answer: B}
		\end{exampleblock}
		\begin{block}{Explanation}
			Demographic parity = same predicted positive rate across groups.
		\end{block}
	\end{frame}
	
	\begin{frame}{Q44 Theory --- Fairness Metrics}
		\begin{block}{Theory}
			\textbf{Group fairness:}
			\begin{itemize}
				\item Demographic parity: equal prediction rates
				\item Equal opportunity: equal true positive rates
				\item Equalized odds: equal TPR and FPR
				\item Predictive parity: equal precision
			\end{itemize}
			\textbf{Individual fairness:} similar individuals get similar predictions.
			\vspace{2mm}
			
			\textbf{Exam tip:} Demographic parity = equal prediction rates.
		\end{block}
	\end{frame}
	
	\begin{frame}{Q45 --- Chatbot Errors}
		\begin{block}{Question}
			Chatbot gives incorrect info. Best action?
		\end{block}
		\begin{enumerate}[label=\Alph*.]
			\item Upgrade model
			\item Develop ethics framework
			\item Periodic risk reports
			\item Human verification mechanisms
		\end{enumerate}
		\begin{exampleblock}{Answer: D}
		\end{exampleblock}
		\begin{block}{Explanation}
			Human verification mitigates incorrect info risk.
		\end{block}
	\end{frame}
	
	\begin{frame}{Q45 Theory --- Human-in-the-Loop}
		\begin{block}{Theory}
			\textbf{Human-in-the-loop:}
			\begin{itemize}
				\item Critical for high-stakes decisions
				\item Mitigates AI errors
				\item Required by many regulations
			\end{itemize}
			\textbf{Applications:}
			\begin{itemize}
				\item Chatbots: human review
				\item Medical: doctor confirms
				\item Credit: loan officer reviews
				\item Fraud: analyst investigates
			\end{itemize}
		\end{block}
	\end{frame}
	
	\begin{frame}{Q46 --- Model Risk vs Bias}
		\begin{block}{Question}
			What is the difference between model risk and algorithmic bias?
		\end{block}
		\begin{enumerate}[label=\Alph*.]
			\item Model risk = errors; algorithmic bias = unfair discrimination
			\item Model risk = unfair discrimination; algorithmic bias = errors
			\item Same thing
			\item Model risk for AI; bias for statistics
		\end{enumerate}
		\begin{exampleblock}{Answer: A}
		\end{exampleblock}
		\begin{block}{Explanation}
			Model risk = errors/misuse. Algorithmic bias = unfair discrimination.
		\end{block}
	\end{frame}
	
	\begin{frame}{Q46 Theory --- Model Risk vs Bias}
		\begin{block}{Theory}
			\textbf{Model risk:}
			\begin{itemize}
				\item Adverse consequences from errors
				\item Sources: assumptions, data, implementation, misuse
				\item Applies to all models
			\end{itemize}
			\textbf{Algorithmic bias:}
			\begin{itemize}
				\item Systematic unfair discrimination
				\item Sources: biased data/features/objectives
				\item Affects protected groups
			\end{itemize}
			\textbf{Relationship:} Bias is a source of model risk.
		\end{block}
	\end{frame}
	
	\begin{frame}{Q47 --- Explainability}
		\begin{block}{Question}
			What is the purpose of explainability?
		\end{block}
		\begin{enumerate}[label=\Alph*.]
			\item Increase accuracy
			\item Provide human-understandable reasons for decisions
			\item Reduce training time
			\item Increase features
		\end{enumerate}
		\begin{exampleblock}{Answer: B}
		\end{exampleblock}
		\begin{block}{Explanation}
			Explainability = human-understandable reasons for decisions.
		\end{block}
	\end{frame}
	
	\begin{frame}{Q47 Theory --- Explainability}
		\begin{block}{Theory}
			\textbf{Explainability:}
			\begin{itemize}
				\item Why did the model decide this?
				\item Which features mattered?
				\item How would predictions change?
			\end{itemize}
			\textbf{Techniques:}
			\begin{itemize}
				\item SHAP: marginal contribution
				\item LIME: local approximation
				\item Feature importance: global ranking
				\item Partial dependence: feature effect
			\end{itemize}
			\textbf{Why:} Compliance (GDPR), trust, validation.
		\end{block}
	\end{frame}
	
	\begin{frame}{Q48 --- Interpretability vs Explainability}
		\begin{block}{Question}
			What is the difference between interpretability and explainability?
		\end{block}
		\begin{enumerate}[label=\Alph*.]
			\item Interpretability = internal mechanisms; explainability = decisions
			\item Interpretability = decisions; explainability = internal
			\item Same thing
			\item Interpretability for linear; explainability for NN
		\end{enumerate}
		\begin{exampleblock}{Answer: A}
		\end{exampleblock}
		\begin{block}{Explanation}
			Interpretability = understand internal workings. Explainability = explain decisions.
		\end{block}
	\end{frame}
	
	\begin{frame}{Q48 Theory --- Interpretability vs Explainability}
		\begin{block}{Theory}
			\textbf{Interpretability:}
			\begin{itemize}
				\item Can you understand internal logic?
				\item Linear regression: yes
				\item Deep NN: no (black box)
			\end{itemize}
			\textbf{Explainability:}
			\begin{itemize}
				\item Can you explain a decision?
				\item Post-hoc for black-box
				\item SHAP, LIME
			\end{itemize}
			\textbf{Exam tip:} Interpretability = internal. Explainability = external explanation.
		\end{block}
	\end{frame}
	
	\begin{frame}{Q49 --- Black Box}
		\begin{block}{Question}
			What is the ``black box'' problem?
		\end{block}
		\begin{enumerate}[label=\Alph*.]
			\item Models stored in black boxes
			\item Decisions cannot be understood by humans
			\item Too large for memory
			\item Trained on black market data
		\end{enumerate}
		\begin{exampleblock}{Answer: B}
		\end{exampleblock}
		\begin{block}{Explanation}
			Black box = decisions cannot be understood.
		\end{block}
	\end{frame}
	
	\begin{frame}{Q49 Theory --- Black Box Problem}
		\begin{block}{Theory}
			\textbf{Black box:}
			\begin{itemize}
				\item Millions of parameters
				\item No simple formula
				\item Cannot trace inputs to outputs
			\end{itemize}
			\textbf{Consequences:}
			\begin{itemize}
				\item Hard to validate
				\item Hard to audit
				\item Hard to trust
				\item Regulatory challenges
			\end{itemize}
			\textbf{Solutions:}
			\begin{itemize}
				\item Interpretable models
				\item Post-hoc explainability
				\item Attention mechanisms
				\item Hybrid models
			\end{itemize}
		\end{block}
	\end{frame}
	
	\begin{frame}{Q50 --- AI in Hiring}
		\begin{block}{Question}
			What is the primary concern with AI in hiring?
		\end{block}
		\begin{enumerate}[label=\Alph*.]
			\item Too slow
			\item May perpetuate historical bias
			\item Too expensive
			\item Cannot process resumes
		\end{enumerate}
		\begin{exampleblock}{Answer: B}
		\end{exampleblock}
		\begin{block}{Explanation}
			AI in hiring may perpetuate historical bias.
		\end{block}
	\end{frame}
	
	\begin{frame}{Q50 Theory --- AI in Hiring}
		\begin{block}{Theory}
			\textbf{Risks:}
			\begin{itemize}
				\item Historical bias
				\item Proxy discrimination
				\item Lack of transparency
				\item Regulatory risk
			\end{itemize}
			\textbf{Mitigation:}
			\begin{itemize}
				\item Audit training data
				\item Remove/de-bias proxies
				\item Fairness constraints
				\item Explanations to candidates
				\item Human review
			\end{itemize}
		\end{block}
	\end{frame}
	
	% ============================================================
	\section{Module 4: Responsible and Ethical AI}
	% ============================================================
	
	\begin{frame}{Q51 --- Shapley Values}
		\begin{block}{Question}
			Which technique explains marginal effect of each feature?
		\end{block}
		\begin{enumerate}[label=\Alph*.]
			\item Shapley values
			\item Feature importance
			\item Grid search
			\item Feature scaling
		\end{enumerate}
		\begin{exampleblock}{Answer: A}
		\end{exampleblock}
		\begin{block}{Explanation}
			Shapley values = marginal contribution of each feature.
		\end{block}
	\end{frame}
	
	\begin{frame}{Q51 Theory --- Shapley Values}
		\begin{block}{Theory}
			\textbf{Shapley values:}
			\begin{itemize}
				\item From game theory
				\item Fairly distributes prediction among features
				\item Considers all subsets
			\end{itemize}
			\textbf{Properties:}
			\begin{itemize}
				\item Additivity
				\item Consistency
				\item Symmetry
			\end{itemize}
			\textbf{vs. Feature importance:}
			\begin{itemize}
				\item Feature importance: global
				\item Shapley: per-prediction
			\end{itemize}
		\end{block}
	\end{frame}
	
	\begin{frame}{Q52 --- Loan Rejection}
		\begin{block}{Question}
			Customers concerned about loan decisions. Source of reputational risk and mitigation?
		\end{block}
		\begin{enumerate}[label=\Alph*.]
			\item Privacy; secure data
			\item Bias; diverse training
			\item Lack of explainability; provide reasons
			\item Lack of transparency; disclose data
		\end{enumerate}
		\begin{exampleblock}{Answer: C}
		\end{exampleblock}
		\begin{block}{Explanation}
			Lack of explainability; provide reasons for rejection.
		\end{block}
	\end{frame}
	
	\begin{frame}{Q52 Theory --- Reputational Risk}
		\begin{block}{Theory}
			\textbf{Sources:}
			\begin{itemize}
				\item Lack of explainability
				\item Algorithmic bias
				\item Privacy breaches
				\item Lack of transparency
			\end{itemize}
			\textbf{Mitigation:}
			\begin{itemize}
				\item Provide reasons
				\item Bias audits
				\item Secure data
				\item Disclose AI use
			\end{itemize}
		\end{block}
	\end{frame}
	
	\begin{frame}{Q53 --- Ethics Framework Advantage}
		\begin{block}{Question}
			Advantage of a practical ethics framework?
		\end{block}
		\begin{enumerate}[label=\Alph*.]
			\item Clear guidelines for ethical dilemmas
			\item Eliminates conflicts of interest
			\item Guarantees legal compliance
			\item Focuses on profit
		\end{enumerate}
		\begin{exampleblock}{Answer: A}
		\end{exampleblock}
		\begin{block}{Explanation}
			Provides clear guidelines for ethical dilemmas.
		\end{block}
	\end{frame}
	
	\begin{frame}{Q53 Theory --- Practical Ethics Frameworks}
		\begin{block}{Theory}
			\textbf{Practical ethics framework:}
			\begin{itemize}
				\item Guidelines for dilemmas
				\item Consistency
				\item Accountability
				\item Builds trust
			\end{itemize}
			\textbf{Does NOT:}
			\begin{itemize}
				\item Eliminate conflicts
				\item Guarantee compliance
				\item Focus on profit
				\item Replace judgment
			\end{itemize}
			\textbf{Components:} Principles, processes, roles.
		\end{block}
	\end{frame}
	
	\begin{frame}{Q54 --- AI Ethics Framework}
		\begin{block}{Question}
			Most likely advantage of AI ethics framework?
		\end{block}
		\begin{enumerate}[label=\Alph*.]
			\item Increases accuracy
			\item Strengthens firewall
			\item Increases testing productivity
			\item Reduces regulatory non-compliance risk
		\end{enumerate}
		\begin{exampleblock}{Answer: D}
		\end{exampleblock}
		\begin{block}{Explanation}
			Reduces regulatory non-compliance risk.
		\end{block}
	\end{frame}
	
	\begin{frame}{Q54 Theory --- AI Ethics and Regulation}
		\begin{block}{Theory}
			\textbf{Why ethics frameworks matter:}
			\begin{itemize}
				\item Regulatory compliance (EU AI Act, GDPR, CCPA)
				\item Reputation
				\item Trust
				\item Risk management
			\end{itemize}
			\textbf{Key regulations:}
			\begin{itemize}
				\item EU AI Act: risk-based
				\item GDPR: right to explanation
				\item CCPA: California privacy
				\item SEC/FCA: model risk
			\end{itemize}
		\end{block}
	\end{frame}
	
	\begin{frame}{Q55 --- Utilitarianism}
		\begin{block}{Question}
			Which is correct regarding Utilitarianism?
		\end{block}
		\begin{enumerate}[label=\Alph*.]
			\item Can lead to violations of individual rights
			\item Focuses on virtuous character
			\item Relies on laws
			\item Values intrinsic morality
		\end{enumerate}
		\begin{exampleblock}{Answer: A}
		\end{exampleblock}
		\begin{block}{Explanation}
			Utilitarianism maximizes overall good; can violate individual rights.
		\end{block}
	\end{frame}
	
	\begin{frame}{Q55 Theory --- Ethical Frameworks}
		\begin{block}{Theory}
			\textbf{Three frameworks:}
			\begin{itemize}
				\item \textbf{Consequentialism (Utilitarianism):} judges by consequences; maximize good; can violate rights
				\item \textbf{Deontology:} judges by rules/duties; inherent right/wrong
				\item \textbf{Virtue Ethics:} character traits; what would a virtuous person do?
			\end{itemize}
			\textbf{Exam tip:} Utilitarianism = consequences. Deontology = rules. Virtue = character.
		\end{block}
	\end{frame}
	
	\begin{frame}{Q56 --- Deontology}
		\begin{block}{Question}
			Which best describes deontology?
		\end{block}
		\begin{enumerate}[label=\Alph*.]
			\item Judges by consequences
			\item Judges by adherence to duties and rules
			\item Single quantifiable metric
			\item Emphasizes virtuous character
		\end{enumerate}
		\begin{exampleblock}{Answer: B}
		\end{exampleblock}
		\begin{block}{Explanation}
			Deontology judges by duties and rules, not consequences.
		\end{block}
	\end{frame}
	
	\begin{frame}{Q56 Theory --- Deontology}
		\begin{block}{Theory}
			\textbf{Deontology:}
			\begin{itemize}
				\item Actions right/wrong in themselves
				\item Consequences do not determine morality
				\item Categorical imperative (Kant)
			\end{itemize}
			\textbf{vs Consequentialism:}
			\begin{itemize}
				\item Consequentialism: ends justify means
				\item Deontology: means matter
			\end{itemize}
			\textbf{AI:} Certain uses inherently wrong (mass surveillance).
		\end{block}
	\end{frame}
	
	\begin{frame}{Q57 --- Justice}
		\begin{block}{Question}
			Which is correct regarding justice in AI ethics?
		\end{block}
		\begin{enumerate}[label=\Alph*.]
			\item Conflicts exist between distributive justice concepts
			\item Impartial application not required for procedural justice
			\item Personal autonomy secondary to social justice
			\item Meritocracy not a distributive justice principle
		\end{enumerate}
		\begin{exampleblock}{Answer: A}
		\end{exampleblock}
		\begin{block}{Explanation}
			Different justice concepts (egalitarianism, utilitarianism) can conflict.
		\end{block}
	\end{frame}
	
	\begin{frame}{Q57 Theory --- Justice}
		\begin{block}{Theory}
			\textbf{Types:}
			\begin{itemize}
				\item Distributive: fair allocation
				\begin{itemize}
					\item Egalitarianism: equal
					\item Utilitarianism: maximize welfare
					\item Meritocracy: based on merit
				\end{itemize}
				\item Procedural: fair processes, impartial law
				\item Social: fairness in society
			\end{itemize}
			\textbf{Conflicts:} Concepts can conflict.
		\end{block}
	\end{frame}
	
	\begin{frame}{Q58 --- Autonomy}
		\begin{block}{Question}
			AI recommends treatment conflicting with patient preferences. Which principle violated?
		\end{block}
		\begin{enumerate}[label=\Alph*.]
			\item Explainability
			\item Beneficence
			\item Autonomy
			\item Non-maleficence
		\end{enumerate}
		\begin{exampleblock}{Answer: C}
		\end{exampleblock}
		\begin{block}{Explanation}
			Autonomy = respect patient's right to decide.
		\end{block}
	\end{frame}
	
	\begin{frame}{Q58 Theory --- AI Ethics Principles}
		\begin{block}{Theory}
			\textbf{Five principles:}
			\begin{itemize}
				\item Nonmaleficence: do no harm
				\item Beneficence: do good
				\item Justice: fairness
				\item Autonomy: respect choice
				\item Explainability: understandable reasons
			\end{itemize}
			\textbf{Autonomy in AI:}
			\begin{itemize}
				\item Right to refuse treatment
				\item Consent to AI decisions
				\item AI supports, not replaces, judgment
			\end{itemize}
		\end{block}
	\end{frame}
	
	\begin{frame}{Q59 --- Sampling Bias}
		\begin{block}{Question}
			Model trained on US travelers, fails in Brazil/Nigeria/Thailand. Cause?
		\end{block}
		\begin{enumerate}[label=\Alph*.]
			\item Problem specification bias
			\item Deployment bias
			\item Historical bias
			\item Sampling bias
		\end{enumerate}
		\begin{exampleblock}{Answer: D}
		\end{exampleblock}
		\begin{block}{Explanation}
			Training sample not representative = sampling bias.
		\end{block}
	\end{frame}
	
	\begin{frame}{Q59 Theory --- Sampling Bias}
		\begin{block}{Theory}
			\textbf{Sampling bias:}
			\begin{itemize}
				\item Training from one population
				\item Deployed in different population
				\item Poor performance on new population
			\end{itemize}
			\textbf{Mitigation:}
			\begin{itemize}
				\item Representative data
				\item Diverse testing
				\item Monitor subgroups
				\item Transfer learning
			\end{itemize}
		\end{block}
	\end{frame}
	
	\begin{frame}{Q60 --- AI Ethics Framework}
		\begin{block}{Question}
			Primary purpose of AI ethics framework?
		\end{block}
		\begin{enumerate}[label=\Alph*.]
			\item Maximize accuracy
			\item Provide guidelines for ethical AI
			\item Reduce training time
			\item Increase features
		\end{enumerate}
		\begin{exampleblock}{Answer: B}
		\end{exampleblock}
		\begin{block}{Explanation}
			Provides guidelines for ethical AI development and use.
		\end{block}
	\end{frame}
	
	\begin{frame}{Q60 Theory --- AI Ethics Frameworks}
		\begin{block}{Theory}
			\textbf{Framework:}
			\begin{itemize}
				\item Defines principles
				\item Establishes processes
				\item Assigns roles
				\item Provides training
			\end{itemize}
			\textbf{Benefits:}
			\begin{itemize}
				\item Reduces regulatory risk
				\item Builds trust
				\item Consistency
				\item Protects reputation
			\end{itemize}
		\end{block}
	\end{frame}
	
	% ============================================================
	\section{Module 5: Data and AI Model Governance}
	% ============================================================
	
	\begin{frame}{Q61 --- Model Landscape}
		\begin{block}{Question}
			Crucial consideration when evaluating model landscape?
		\end{block}
		\begin{enumerate}[label=\Alph*.]
			\item Models operate independently
			\item All use same input data
			\item Consolidate models
			\item Interdependencies identified
		\end{enumerate}
		\begin{exampleblock}{Answer: D}
		\end{exampleblock}
		\begin{block}{Explanation}
			Model landscape shows interactions; identify interdependencies.
		\end{block}
	\end{frame}
	
	\begin{frame}{Q61 Theory --- Model Landscape}
		\begin{block}{Theory}
			\textbf{Model landscape:}
			\begin{itemize}
				\item Shows relationships
				\item Identifies critical models
				\item Maps data flows
				\item Highlights interdependencies
			\end{itemize}
			\textbf{Why:}
			\begin{itemize}
				\item One output = another's input
				\item Failure cascades
				\item Correlated errors
				\item Systemic risk
			\end{itemize}
		\end{block}
	\end{frame}
	
	\begin{frame}{Q62 --- Model Governance Roles}
		\begin{block}{Question}
			Which correctly describes roles?
		\end{block}
		\begin{enumerate}[label=\Alph*.]
			\item Board maintains inventory
			\item Model risk function sets appetite
			\item Model ownership never shared
			\item Auditors are independent checks
		\end{enumerate}
		\begin{exampleblock}{Answer: D}
		\end{exampleblock}
		\begin{block}{Explanation}
			Auditors check validation. Board sets appetite. Model risk function maintains inventory.
		\end{block}
	\end{frame}
	
	\begin{frame}{Q62 Theory --- Model Governance Roles}
		\begin{block}{Theory}
			\textbf{Roles:}
			\begin{itemize}
				\item Board: sets risk appetite
				\item Model risk function: inventory, policies
				\item Developers: build models
				\item Validators: independent validation
				\item Internal auditors: check validation
				\item External auditors: external check
				\item Model owners: accountable
			\end{itemize}
		\end{block}
	\end{frame}
	
	\begin{frame}{Q63 --- IT Integration}
		\begin{block}{Question}
			Which team ensures IT integration without disruption?
		\end{block}
		\begin{enumerate}[label=\Alph*.]
			\item Core Implementation
			\item System Implementation
			\item Model Design
			\item Data Management
		\end{enumerate}
		\begin{exampleblock}{Answer: B}
		\end{exampleblock}
		\begin{block}{Explanation}
			System Implementation Team integrates model with IT systems.
		\end{block}
	\end{frame}
	
	\begin{frame}{Q63 Theory --- Implementation Teams}
		\begin{block}{Theory}
			\textbf{Teams:}
			\begin{itemize}
				\item Model Design: architecture, methodology
				\item Core Implementation: builds/codes
				\item System Implementation: IT integration
				\item Data Management: data quality
			\end{itemize}
			\textbf{System Implementation:}
			\begin{itemize}
				\item IT integration
				\item Production function
				\item Deployment
				\item Monitoring
			\end{itemize}
		\end{block}
	\end{frame}
	
	\begin{frame}{Q64 --- Next Step After Testing}
		\begin{block}{Question}
			New model outperforms traditional. Next step?
		\end{block}
		\begin{enumerate}[label=\Alph*.]
			\item Deploy to production
			\item Inform IT
			\item Retire existing model
			\item Request validation
		\end{enumerate}
		\begin{exampleblock}{Answer: D}
		\end{exampleblock}
		\begin{block}{Explanation}
			Validate before deployment. Never deploy without validation.
		\end{block}
	\end{frame}
	
	\begin{frame}{Q64 Theory --- Model Validation}
		\begin{block}{Theory}
			\textbf{Validation:}
			\begin{itemize}
				\item Independent team
				\item Conceptual soundness
				\item Data quality
				\item Performance testing
				\item Implementation review
			\end{itemize}
			\textbf{When:}
			\begin{itemize}
				\item Before deployment
				\item Periodically
				\item After changes
				\item On degradation
			\end{itemize}
		\end{block}
	\end{frame}
	
	\begin{frame}{Q65 --- Alternative Data}
		\begin{block}{Question}
			AI credit model uses school transcripts, social media. Validation finding?
		\end{block}
		\begin{enumerate}[label=\Alph*.]
			\item No concerns
			\item Raise concerns about sensitive data
			\item Approve if backtest improves
			\item Add more alternative data
		\end{enumerate}
		\begin{exampleblock}{Answer: B}
		\end{exampleblock}
		\begin{block}{Explanation}
			Sensitive alternative data requires heightened scrutiny.
		\end{block}
	\end{frame}
	
	\begin{frame}{Q65 Theory --- Alternative Data}
		\begin{block}{Theory}
			\textbf{Alternative data:}
			\begin{itemize}
				\item Social media, transcripts, rent
				\item Expands credit access
				\item Ethical/legal concerns
			\end{itemize}
			\textbf{Governance:}
			\begin{itemize}
				\item Provenance: legal right
				\item Consent: applicant approval
				\item Privacy: protect info
				\item Fairness: avoid discrimination
				\item Transparency: disclose use
			\end{itemize}
		\end{block}
	\end{frame}
	
	\begin{frame}{Q66 --- Model Monitoring}
		\begin{block}{Question}
			Advantage of increasing monitoring frequency?
		\end{block}
		\begin{enumerate}[label=\Alph*.]
			\item Reduces need for retraining
			\item Detects degradation or drift
			\item Minimizes IT involvement
			\item Improves interpretability
		\end{enumerate}
		\begin{exampleblock}{Answer: B}
		\end{exampleblock}
		\begin{block}{Explanation}
			Detects degradation/drift for timely intervention.
		\end{block}
	\end{frame}
	
	\begin{frame}{Q66 Theory --- Monitoring and Drift}
		\begin{block}{Theory}
			\textbf{Monitoring:}
			\begin{itemize}
				\item Detect performance degradation
				\item Identify data drift
				\item Identify concept drift
				\item Trigger retraining
			\end{itemize}
			\textbf{Drift types:}
			\begin{itemize}
				\item Data: input distribution changes
				\item Concept: input-output relationship changes
				\item Label: target distribution changes
			\end{itemize}
		\end{block}
	\end{frame}
	
	\begin{frame}{Q67 --- Model Inventory}
		\begin{block}{Question}
			Purpose of a model inventory?
		\end{block}
		\begin{enumerate}[label=\Alph*.]
			\item Store code
			\item List all models and track risk
			\item Train models
			\item Visualize performance
		\end{enumerate}
		\begin{exampleblock}{Answer: B}
		\end{exampleblock}
		\begin{block}{Explanation}
			Lists all models with purpose, risk, owner, validation status.
		\end{block}
	\end{frame}
	
	\begin{frame}{Q67 Theory --- Model Inventory}
		\begin{block}{Theory}
			\textbf{Model inventory:}
			\begin{itemize}
				\item Maintained by model risk function
				\item Tracks all models
				\item Documents purpose, owner, risk
				\item Records validation status
				\item Identifies interdependencies
			\end{itemize}
			\textbf{Contents:}
			\begin{itemize}
				\item Name, description
				\item Purpose
				\item Owner, developer
				\item Risk classification
				\item Validation status
			\end{itemize}
		\end{block}
	\end{frame}
	
	\begin{frame}{Q68 --- Model Documentation}
		\begin{block}{Question}
			Purpose of model documentation?
		\end{block}
		\begin{enumerate}[label=\Alph*.]
			\item Increase accuracy
			\item Record design, assumptions, limitations
			\item Speed up training
			\item Reduce features
		\end{enumerate}
		\begin{exampleblock}{Answer: B}
		\end{exampleblock}
		\begin{block}{Explanation}
			Records design, assumptions, limitations for validation and audit.
		\end{block}
	\end{frame}
	
	\begin{frame}{Q68 Theory --- Model Documentation}
		\begin{block}{Theory}
			\textbf{Documentation:}
			\begin{itemize}
				\item Purpose and use
				\item Methodology
				\item Data sources
				\item Limitations
				\item Performance
				\item Validation results
				\item Usage guidelines
			\end{itemize}
			\textbf{Why:} Validation, audit, knowledge transfer, compliance.
		\end{block}
	\end{frame}
	
	\begin{frame}{Q69 --- Use Test}
		\begin{block}{Question}
			Purpose of use test?
		\end{block}
		\begin{enumerate}[label=\Alph*.]
			\item Test production use
			\item Assess appropriateness for intended use
			\item Test speed
			\item Test accuracy
		\end{enumerate}
		\begin{exampleblock}{Answer: B}
		\end{exampleblock}
		\begin{block}{Explanation}
			Use test assesses fit for intended purpose.
		\end{block}
	\end{frame}
	
	\begin{frame}{Q69 Theory --- Use Test}
		\begin{block}{Theory}
			\textbf{Use test:}
			\begin{itemize}
				\item Used as intended?
				\item Users understand limitations?
				\item Appropriate for context?
				\item Outputs interpreted correctly?
			\end{itemize}
			\textbf{Questions:}
			\begin{itemize}
				\item What decisions supported?
				\item Who uses it?
				\item Are assumptions understood?
				\item Alternatives?
			\end{itemize}
		\end{block}
	\end{frame}
	
	\begin{frame}{Q70 --- Testing}
		\begin{block}{Question}
			Difference between white-box and black-box testing?
		\end{block}
		\begin{enumerate}[label=\Alph*.]
			\item White-box = internal logic; black-box = inputs/outputs
			\item White-box = inputs/outputs; black-box = internal
			\item Same
			\item White-box for AI; black-box for statistics
		\end{enumerate}
		\begin{exampleblock}{Answer: A}
		\end{exampleblock}
		\begin{block}{Explanation}
			White-box = internal. Black-box = external.
		\end{block}
	\end{frame}
	
	\begin{frame}{Q70 Theory --- Testing}
		\begin{block}{Theory}
			\textbf{White-box:}
			\begin{itemize}
				\item Internal logic and code
				\item Model structure knowledge
				\item Component tests
				\item Implementation errors
			\end{itemize}
			\textbf{Black-box:}
			\begin{itemize}
				\item Inputs/outputs only
				\item No internal knowledge
				\item Overall behavior
				\item Real-world simulation
			\end{itemize}
			\textbf{Both needed.}
		\end{block}
	\end{frame}
	
	\begin{frame}{Q71 --- Data Provenance}
		\begin{block}{Question}
			Purpose of data provenance?
		\end{block}
		\begin{enumerate}[label=\Alph*.]
			\item Verify legal right to use data
			\item Increase volume
			\item Speed up training
			\item Reduce quality
		\end{enumerate}
		\begin{exampleblock}{Answer: A}
		\end{exampleblock}
		\begin{block}{Explanation}
			Provenance = legal right to use data.
		\end{block}
	\end{frame}
	
	\begin{frame}{Q71 Theory --- Data Provenance}
		\begin{block}{Theory}
			\textbf{Data provenance:}
			\begin{itemize}
				\item Where from?
				\item Legal right?
				\item Consent obtained?
				\item Accurate?
			\end{itemize}
			\textbf{Why:}
			\begin{itemize}
				\item Legal compliance (GDPR, CCPA)
				\item Ethical use
				\item Regulatory
				\item Reputation
			\end{itemize}
		\end{block}
	\end{frame}
	
	\begin{frame}{Q72 --- Data Classification}
		\begin{block}{Question}
			Purpose of data classification?
		\end{block}
		\begin{enumerate}[label=\Alph*.]
			\item Increase volume
			\item Categorize by sensitivity and structure
			\item Speed up training
			\item Reduce features
		\end{enumerate}
		\begin{exampleblock}{Answer: B}
		\end{exampleblock}
		\begin{block}{Explanation}
			Categorizes by sensitivity (public, internal, confidential) and structure.
		\end{block}
	\end{frame}
	
	\begin{frame}{Q72 Theory --- Data Classification}
		\begin{block}{Theory}
			\textbf{Classification:}
			\begin{itemize}
				\item Sensitivity:
				\begin{itemize}
					\item Public, Internal, Confidential, Restricted
				\end{itemize}
				\item Structure:
				\begin{itemize}
					\item Structured (tables), Unstructured (text, images)
				\end{itemize}
			\end{itemize}
			\textbf{Purpose:} Determine handling, access, protection.
		\end{block}
	\end{frame}
	
	\begin{frame}{Q73 --- Metadata}
		\begin{block}{Question}
			Purpose of metadata management?
		\end{block}
		\begin{enumerate}[label=\Alph*.]
			\item Increase volume
			\item Document origin, structure, definitions, relationships
			\item Speed up training
			\item Reduce quality
		\end{enumerate}
		\begin{exampleblock}{Answer: B}
		\end{exampleblock}
		\begin{block}{Explanation}
			Metadata = data about data.
		\end{block}
	\end{frame}
	
	\begin{frame}{Q73 Theory --- Metadata Management}
		\begin{block}{Theory}
			\textbf{Metadata:}
			\begin{itemize}
				\item Origin, history
				\item Structure, format
				\item Definitions, meaning
				\item Relationships
				\item Access permissions
			\end{itemize}
			\textbf{Why:}
			\begin{itemize}
				\item Discovery
				\item Lineage
				\item Compliance
				\item Governance
			\end{itemize}
		\end{block}
	\end{frame}
	
	\begin{frame}{Q74 --- Data Access}
		\begin{block}{Question}
			Purpose of data access controls?
		\end{block}
		\begin{enumerate}[label=\Alph*.]
			\item Increase volume
			\item Only authorized users access data
			\item Speed up training
			\item Reduce quality
		\end{enumerate}
		\begin{exampleblock}{Answer: B}
		\end{exampleblock}
		\begin{block}{Explanation}
			Ensures only authorized users access data.
		\end{block}
	\end{frame}
	
	\begin{frame}{Q74 Theory --- Data Access}
		\begin{block}{Theory}
			\textbf{Access controls:}
			\begin{itemize}
				\item RBAC: role-based
				\item ABAC: attribute-based
				\item Least privilege
				\item Need-to-know
			\end{itemize}
			\textbf{Why:}
			\begin{itemize}
				\item Privacy
				\item Security
				\item Compliance
				\item Prevent misuse
			\end{itemize}
		\end{block}
	\end{frame}
	
	\begin{frame}{Q75 --- Data Governance}
		\begin{block}{Question}
			Purpose of data governance framework?
		\end{block}
		\begin{enumerate}[label=\Alph*.]
			\item Increase volume
			\item Ensure quality, security, compliance
			\item Speed up training
			\item Reduce features
		\end{enumerate}
		\begin{exampleblock}{Answer: B}
		\end{exampleblock}
		\begin{block}{Explanation}
			Ensures data quality, security, compliance, proper use.
		\end{block}
	\end{frame}
	
	\begin{frame}{Q75 Theory --- Data Governance}
		\begin{block}{Theory}
			\textbf{Data governance:}
			\begin{itemize}
				\item Strategy and vision
				\item Quality
				\item Provenance
				\item Classification
				\item Metadata
				\item Protection
				\item Access
				\item Compliance
				\item Roles
			\end{itemize}
		\end{block}
	\end{frame}
	
	\begin{frame}{Q76 --- Model Risk Management}
		\begin{block}{Question}
			Purpose of model risk management?
		\end{block}
		\begin{enumerate}[label=\Alph*.]
			\item Increase accuracy
			\item Identify, assess, mitigate model risks
			\item Speed up training
			\item Reduce features
		\end{enumerate}
		\begin{exampleblock}{Answer: B}
		\end{exampleblock}
		\begin{block}{Explanation}
			Identifies, assesses, mitigates model risks.
		\end{block}
	\end{frame}
	
	\begin{frame}{Q76 Theory --- Model Risk Management}
		\begin{block}{Theory}
			\textbf{Model risk management:}
			\begin{itemize}
				\item Identification, inventory
				\item Risk assessment
				\item Development, documentation
				\item Independent validation
				\item Monitoring
				\item Governance
			\end{itemize}
			\textbf{Sources:}
			\begin{itemize}
				\item Wrong assumptions
				\item Bad data
				\item Implementation errors
				\item Misuse
				\item Drift
			\end{itemize}
		\end{block}
	\end{frame}
	
	\begin{frame}{Q77 --- Model Risk Function}
		\begin{block}{Question}
			Role of model risk function?
		\end{block}
		\begin{enumerate}[label=\Alph*.]
			\item Build models
			\item Maintain inventory and develop policies
			\item Use models
			\item Reduce features
		\end{enumerate}
		\begin{exampleblock}{Answer: B}
		\end{exampleblock}
		\begin{block}{Explanation}
			Maintains inventory, develops policies, oversees model risk.
		\end{block}
	\end{frame}
	
	\begin{frame}{Q77 Theory --- Model Risk Function}
		\begin{block}{Theory}
			\textbf{Responsibilities:}
			\begin{itemize}
				\item Maintain inventory
				\item Develop policies
				\item Classify risk
				\item Coordinate validation
				\item Monitor performance
				\item Report to management
				\item Ensure compliance
			\end{itemize}
			\textbf{Independence:} From developers.
		\end{block}
	\end{frame}
	
	\begin{frame}{Q78 --- Periodic Validation}
		\begin{block}{Question}
			Purpose of periodic model validation?
		\end{block}
		\begin{enumerate}[label=\Alph*.]
			\item Increase accuracy
			\item Confirm model remains fit for purpose
			\item Speed up training
			\item Reduce features
		\end{enumerate}
		\begin{exampleblock}{Answer: B}
		\end{exampleblock}
		\begin{block}{Explanation}
			Confirms model remains fit for purpose over time.
		\end{block}
	\end{frame}
	
	\begin{frame}{Q78 Theory --- Periodic Validation}
		\begin{block}{Theory}
			\textbf{Periodic validation:}
			\begin{itemize}
				\item Assumptions valid
				\item Data appropriate
				\item Performance as expected
				\item Methodology current
				\item Aligned with purpose
			\end{itemize}
			\textbf{Frequency:}
			\begin{itemize}
				\item Based on risk ranking
				\item Higher risk = more frequent
				\item Annually or risk-based
				\item Triggered by changes
			\end{itemize}
		\end{block}
	\end{frame}
	
	\begin{frame}{Q79 --- Implementation}
		\begin{block}{Question}
			Purpose of model implementation tasks?
		\end{block}
		\begin{enumerate}[label=\Alph*.]
			\item Design model
			\item Deploy to production and integrate IT
			\item Validate model
			\item Reduce features
		\end{enumerate}
		\begin{exampleblock}{Answer: B}
		\end{exampleblock}
		\begin{block}{Explanation}
			Deploy to production, integrate IT systems.
		\end{block}
	\end{frame}
	
	\begin{frame}{Q79 Theory --- Model Implementation}
		\begin{block}{Theory}
			\textbf{Implementation:}
			\begin{itemize}
				\item IT integration
				\item Data pipeline
				\item Performance testing
				\item User training
				\item Documentation
				\item Monitoring setup
			\end{itemize}
			\textbf{Challenges:}
			\begin{itemize}
				\item Legacy systems
				\item Data quality
				\item Scalability
				\item Security
			\end{itemize}
		\end{block}
	\end{frame}
	
	\begin{frame}{Q80 --- Adaptation}
		\begin{block}{Question}
			Purpose of model adaptation tasks?
		\end{block}
		\begin{enumerate}[label=\Alph*.]
			\item Design model
			\item Update model for changing conditions
			\item Validate model
			\item Reduce features
		\end{enumerate}
		\begin{exampleblock}{Answer: B}
		\end{exampleblock}
		\begin{block}{Explanation}
			Updates model for changing conditions.
		\end{block}
	\end{frame}
	
	\begin{frame}{Q80 Theory --- Model Adaptation}
		\begin{block}{Theory}
			\textbf{Adaptation:}
			\begin{itemize}
				\item Retraining
				\item Recalibration
				\item Update assumptions
				\item Add features
				\item Regulatory response
			\end{itemize}
			\textbf{Triggers:}
			\begin{itemize}
				\item Degradation
				\item Data drift
				\item Concept drift
				\item Business changes
				\item Regulatory changes
			\end{itemize}
		\end{block}
	\end{frame}
	
	\begin{frame}{Q81 --- Risk Appetite}
		\begin{block}{Question}
			Purpose of model risk appetite?
		\end{block}
		\begin{enumerate}[label=\Alph*.]
			\item Increase accuracy
			\item Define acceptable model risk level
			\item Speed up training
			\item Reduce features
		\end{enumerate}
		\begin{exampleblock}{Answer: B}
		\end{exampleblock}
		\begin{block}{Explanation}
			Defines acceptable model risk, set by board.
		\end{block}
	\end{frame}
	
	\begin{frame}{Q81 Theory --- Model Risk Appetite}
		\begin{block}{Theory}
			\textbf{Risk appetite:}
			\begin{itemize}
				\item Set by board
				\item Defines tolerance
				\item Guides management
				\item Aligns with overall appetite
			\end{itemize}
			\textbf{Considerations:}
			\begin{itemize}
				\item Business strategy
				\item Regulatory
				\item Risk capacity
				\item Stakeholder expectations
			\end{itemize}
		\end{block}
	\end{frame}
	
	\begin{frame}{Q82 --- Governance Policies}
		\begin{block}{Question}
			Purpose of model governance policies?
		\end{block}
		\begin{enumerate}[label=\Alph*.]
			\item Increase accuracy
			\item Establish rules for development, validation, use
			\item Speed up training
			\item Reduce features
		\end{enumerate}
		\begin{exampleblock}{Answer: B}
		\end{exampleblock}
		\begin{block}{Explanation}
			Establishes rules for consistency and accountability.
		\end{block}
	\end{frame}
	
	\begin{frame}{Q82 Theory --- Model Governance Policies}
		\begin{block}{Theory}
			\textbf{Policies:}
			\begin{itemize}
				\item Development standards
				\item Validation requirements
				\item Documentation standards
				\item Approval processes
				\item Monitoring
				\item Roles
			\end{itemize}
			\textbf{Purpose:}
			\begin{itemize}
				\item Consistency
				\item Accountability
				\item Compliance
				\item Risk management
			\end{itemize}
		\end{block}
	\end{frame}
	
	\begin{frame}{Q83 --- Documentation Standards}
		\begin{block}{Question}
			Purpose of documentation standards?
		\end{block}
		\begin{enumerate}[label=\Alph*.]
			\item Increase accuracy
			\item Ensure consistent documentation
			\item Speed up training
			\item Reduce features
		\end{enumerate}
		\begin{exampleblock}{Answer: B}
		\end{exampleblock}
		\begin{block}{Explanation}
			Ensures consistent and comprehensive documentation.
		\end{block}
	\end{frame}
	
	\begin{frame}{Q83 Theory --- Documentation Standards}
		\begin{block}{Theory}
			\textbf{Standards:}
			\begin{itemize}
				\item Purpose
				\item Methodology
				\item Data sources
				\item Limitations
				\item Performance
				\item Validation
				\item Usage
			\end{itemize}
			\textbf{Benefits:}
			\begin{itemize}
				\item Consistency
				\item Knowledge transfer
				\item Audit readiness
				\item Compliance
			\end{itemize}
		\end{block}
	\end{frame}
	
	\begin{frame}{Q84 --- Inventory and Landscape}
		\begin{block}{Question}
			Purpose of model inventory and landscape?
		\end{block}
		\begin{enumerate}[label=\Alph*.]
			\item Increase accuracy
			\item Track all models and interdependencies
			\item Speed up training
			\item Reduce features
		\end{enumerate}
		\begin{exampleblock}{Answer: B}
		\end{exampleblock}
		\begin{block}{Explanation}
			Inventory = list. Landscape = map of interactions.
		\end{block}
	\end{frame}
	
	\begin{frame}{Q84 Theory --- Inventory vs Landscape}
		\begin{block}{Theory}
			\textbf{Inventory:}
			\begin{itemize}
				\item List of all models
				\item Name, purpose, owner
				\item Risk level
				\item Validation status
			\end{itemize}
			\textbf{Landscape:}
			\begin{itemize}
				\item Visual interactions
				\item Data flows
				\item Critical models
				\item Interdependencies
			\end{itemize}
		\end{block}
	\end{frame}
	
	\begin{frame}{Q85 --- Model Risk Reporting}
		\begin{block}{Question}
			Purpose of model risk reporting?
		\end{block}
		\begin{enumerate}[label=\Alph*.]
			\item Increase accuracy
			\item Inform management about model risk
			\item Speed up training
			\item Reduce features
		\end{enumerate}
		\begin{exampleblock}{Answer: B}
		\end{exampleblock}
		\begin{block}{Explanation}
			Informs management about model risk.
		\end{block}
	\end{frame}
	
	\begin{frame}{Q85 Theory --- Reporting}
		\begin{block}{Theory}
			\textbf{Reporting:}
			\begin{itemize}
				\item Inventory summary
				\item Validation status
				\item Key issues
				\item Performance
				\item Remediation
			\end{itemize}
			\textbf{Audience:}
			\begin{itemize}
				\item Board
				\item Senior management
				\item Risk committee
				\item Regulators
			\end{itemize}
		\end{block}
	\end{frame}
	
	\begin{frame}{Q86 --- Change Management}
		\begin{block}{Question}
			Purpose of model change management?
		\end{block}
		\begin{enumerate}[label=\Alph*.]
			\item Increase accuracy
			\item Control changes with review
			\item Speed up training
			\item Reduce features
		\end{enumerate}
		\begin{exampleblock}{Answer: B}
		\end{exampleblock}
		\begin{block}{Explanation}
			Controls changes with review, approval, documentation.
		\end{block}
	\end{frame}
	
	\begin{frame}{Q86 Theory --- Change Management}
		\begin{block}{Theory}
			\textbf{Change management:}
			\begin{itemize}
				\item Identify need
				\item Assess impact
				\item Obtain approval
				\item Implement
				\item Document
				\item Validate
			\end{itemize}
			\textbf{Types:}
			\begin{itemize}
				\item Minor: parameter updates
				\item Major: methodology changes
				\item Emergency: urgent fixes
			\end{itemize}
		\end{block}
	\end{frame}
	
	\begin{frame}{Q87 --- Model Review}
		\begin{block}{Question}
			Purpose of model review?
		\end{block}
		\begin{enumerate}[label=\Alph*.]
			\item Increase accuracy
			\item Periodically assess performance and fitness
			\item Speed up training
			\item Reduce features
		\end{enumerate}
		\begin{exampleblock}{Answer: B}
		\end{exampleblock}
		\begin{block}{Explanation}
			Periodically assesses performance and fitness.
		\end{block}
	\end{frame}
	
	\begin{frame}{Q87 Theory --- Model Review}
		\begin{block}{Theory}
			\textbf{Review:}
			\begin{itemize}
				\item Performance
				\item Assumptions
				\item Data quality
				\item Methodology
				\item Documentation
			\end{itemize}
			\textbf{Outcomes:}
			\begin{itemize}
				\item Continue
				\item Modify
				\item Replace
				\item Retire
			\end{itemize}
		\end{block}
	\end{frame}
	
	\begin{frame}{Q88 --- Retirement}
		\begin{block}{Question}
			Purpose of model retirement?
		\end{block}
		\begin{enumerate}[label=\Alph*.]
			\item Increase accuracy
			\item Formally decommission unused model
			\item Speed up training
			\item Reduce features
		\end{enumerate}
		\begin{exampleblock}{Answer: B}
		\end{exampleblock}
		\begin{block}{Explanation}
			Formally decommissions unused models.
		\end{block}
	\end{frame}
	
	\begin{frame}{Q88 Theory --- Retirement}
		\begin{block}{Theory}
			\textbf{Retirement:}
			\begin{itemize}
				\item Document reason
				\item Obtain approval
				\item Transition
				\item Archive
				\item Update inventory
			\end{itemize}
			\textbf{Reasons:}
			\begin{itemize}
				\item No longer needed
				\item Better model
				\item Regulatory changes
				\item Business changes
			\end{itemize}
		\end{block}
	\end{frame}
	
	\begin{frame}{Q89 --- Risk Culture}
		\begin{block}{Question}
			Purpose of model risk culture?
		\end{block}
		\begin{enumerate}[label=\Alph*.]
			\item Increase accuracy
			\item Promote awareness and accountability
			\item Speed up training
			\item Reduce features
		\end{enumerate}
		\begin{exampleblock}{Answer: B}
		\end{exampleblock}
		\begin{block}{Explanation}
			Promotes awareness and accountability for model risk.
		\end{block}
	\end{frame}
	
	\begin{frame}{Q89 Theory --- Risk Culture}
		\begin{block}{Theory}
			\textbf{Culture:}
			\begin{itemize}
				\item Awareness
				\item Accountability
				\item Open communication
				\item Continuous learning
				\item Ethical behavior
			\end{itemize}
			\textbf{Building blocks:}
			\begin{itemize}
				\item Tone from top
				\item Training
				\item Aligned incentives
				\item Transparency
			\end{itemize}
		\end{block}
	\end{frame}
	
	\begin{frame}{Q90 --- Training}
		\begin{block}{Question}
			Purpose of model risk training?
		\end{block}
		\begin{enumerate}[label=\Alph*.]
			\item Increase accuracy
			\item Ensure staff understand model risk
			\item Speed up training
			\item Reduce features
		\end{enumerate}
		\begin{exampleblock}{Answer: B}
		\end{exampleblock}
		\begin{block}{Explanation}
			Ensures staff understand model risk and responsibilities.
		\end{block}
	\end{frame}
	
	\begin{frame}{Q90 Theory --- Training}
		\begin{block}{Theory}
			\textbf{Training:}
			\begin{itemize}
				\item What is model risk?
				\item Roles and responsibilities
				\item Policies
				\item Best practices
				\item Case studies
			\end{itemize}
			\textbf{Audience:}
			\begin{itemize}
				\item Board
				\item Management
				\item Developers, validators
				\item Users, all staff
			\end{itemize}
		\end{block}
	\end{frame}
	
	\begin{frame}{Q91 --- Inscrutability (repeat)}
		\begin{block}{Question}
			Which statement correctly describes inscrutability?
		\end{block}
		\begin{enumerate}[label=\Alph*.]
			\item Cannot understand internal mechanisms
			\item Cannot determine accuracy
			\item Cannot explain predictions
			\item Cannot generalize
		\end{enumerate}
		\begin{exampleblock}{Answer: A}
		\end{exampleblock}
		\begin{block}{Explanation}
			Inscrutability = cannot understand internal mechanisms.
		\end{block}
	\end{frame}
	
	\begin{frame}{Q91 Theory --- Inscrutability in Model Risk}
		\begin{block}{Theory}
			\textbf{Inscrutability:}
			\begin{itemize}
				\item Cannot understand internal workings
				\item Hard to validate, audit
				\item Regulatory challenges
			\end{itemize}
			\textbf{Mitigation:}
			\begin{itemize}
				\item Interpretable models
				\item Post-hoc explainability
				\item Documentation
				\item Human oversight
			\end{itemize}
		\end{block}
	\end{frame}
	
	\begin{frame}{Q92 --- Data Governance}
		\begin{block}{Question}
			Primary purpose of data governance framework?
		\end{block}
		\begin{enumerate}[label=\Alph*.]
			\item Increase volume
			\item Ensure quality, security, compliance
			\item Speed up training
			\item Reduce features
		\end{enumerate}
		\begin{exampleblock}{Answer: B}
		\end{exampleblock}
		\begin{block}{Explanation}
			Ensures quality, security, compliance.
		\end{block}
	\end{frame}
	
	\begin{frame}{Q92 Theory --- Data Governance}
		\begin{block}{Theory}
			\textbf{Data governance:}
			\begin{itemize}
				\item Strategy
				\item Quality
				\item Provenance
				\item Classification
				\item Metadata
				\item Protection
				\item Access
				\item Compliance
				\item Roles
			\end{itemize}
		\end{block}
	\end{frame}
	
	\begin{frame}{Q93 --- Risk Appetite (repeat)}
		\begin{block}{Question}
			Purpose of model risk appetite?
		\end{block}
		\begin{enumerate}[label=\Alph*.]
			\item Increase accuracy
			\item Define acceptable risk level
			\item Speed up training
			\item Reduce features
		\end{enumerate}
		\begin{exampleblock}{Answer: B}
		\end{exampleblock}
		\begin{block}{Explanation}
			Defines acceptable model risk level.
		\end{block}
	\end{frame}
	
	\begin{frame}{Q93 Theory --- Risk Appetite}
		\begin{block}{Theory}
			\textbf{Risk appetite:}
			\begin{itemize}
				\item Set by board
				\item Defines tolerance
				\item Guides management
				\item Aligns with overall appetite
			\end{itemize}
		\end{block}
	\end{frame}
	
	\begin{frame}{Q94 --- Validation (repeat)}
		\begin{block}{Question}
			Purpose of model validation?
		\end{block}
		\begin{enumerate}[label=\Alph*.]
			\item Increase accuracy
			\item Independently assess soundness and fitness
			\item Speed up training
			\item Reduce features
		\end{enumerate}
		\begin{exampleblock}{Answer: B}
		\end{exampleblock}
		\begin{block}{Explanation}
			Independently assesses model soundness and fitness.
		\end{block}
	\end{frame}
	
	\begin{frame}{Q94 Theory --- Validation}
		\begin{block}{Theory}
			\textbf{Validation:}
			\begin{itemize}
				\item Conceptual soundness
				\item Data quality
				\item Performance
				\item Implementation
				\item Documentation
			\end{itemize}
			\textbf{When:}
			\begin{itemize}
				\item Before deployment
				\item Periodically
				\item After changes
			\end{itemize}
		\end{block}
	\end{frame}
	
	\begin{frame}{Q95 --- Monitoring (repeat)}
		\begin{block}{Question}
			Purpose of model monitoring?
		\end{block}
		\begin{enumerate}[label=\Alph*.]
			\item Increase accuracy
			\item Detect degradation and drift
			\item Speed up training
			\item Reduce features
		\end{enumerate}
		\begin{exampleblock}{Answer: B}
		\end{exampleblock}
		\begin{block}{Explanation}
			Detects degradation and drift.
		\end{block}
	\end{frame}
	
	\begin{frame}{Q95 Theory --- Monitoring}
		\begin{block}{Theory}
			\textbf{Monitoring:}
			\begin{itemize}
				\item Performance metrics
				\item Data drift
				\item Concept drift
				\item Alerts
				\item Reporting
			\end{itemize}
			\textbf{Triggers:}
			\begin{itemize}
				\item Below threshold
				\item Significant drift
				\item Regulatory changes
				\item Business changes
			\end{itemize}
		\end{block}
	\end{frame}
	
	\begin{frame}{Q96 --- Documentation (repeat)}
		\begin{block}{Question}
			Purpose of model documentation?
		\end{block}
		\begin{enumerate}[label=\Alph*.]
			\item Increase accuracy
			\item Record design, assumptions, limitations
			\item Speed up training
			\item Reduce features
		\end{enumerate}
		\begin{exampleblock}{Answer: B}
		\end{exampleblock}
		\begin{block}{Explanation}
			Records design, assumptions, limitations.
		\end{block}
	\end{frame}
	
	\begin{frame}{Q96 Theory --- Documentation}
		\begin{block}{Theory}
			\textbf{Documentation:}
			\begin{itemize}
				\item Purpose
				\item Methodology
				\item Data sources
				\item Limitations
				\item Performance
				\item Validation
				\item Usage
			\end{itemize}
		\end{block}
	\end{frame}
	
	\begin{frame}{Q97 --- Inventory (repeat)}
		\begin{block}{Question}
			Purpose of model inventory?
		\end{block}
		\begin{enumerate}[label=\Alph*.]
			\item Increase accuracy
			\item List all models and track risk
			\item Speed up training
			\item Reduce features
		\end{enumerate}
		\begin{exampleblock}{Answer: B}
		\end{exampleblock}
		\begin{block}{Explanation}
			Lists all models with risk and validation status.
		\end{block}
	\end{frame}
	
	\begin{frame}{Q97 Theory --- Inventory}
		\begin{block}{Theory}
			\textbf{Inventory:}
			\begin{itemize}
				\item Name, description
				\item Purpose, use case
				\item Owner, developer
				\item Risk classification
				\item Validation status
				\item Last review
			\end{itemize}
			\textbf{Maintained by:} Model risk function.
		\end{block}
	\end{frame}
	
	\begin{frame}{Q98 --- Landscape (repeat)}
		\begin{block}{Question}
			Purpose of model landscape?
		\end{block}
		\begin{enumerate}[label=\Alph*.]
			\item Increase accuracy
			\item Show interactions and interdependencies
			\item Speed up training
			\item Reduce features
		\end{enumerate}
		\begin{exampleblock}{Answer: B}
		\end{exampleblock}
		\begin{block}{Explanation}
			Shows interactions and interdependencies.
		\end{block}
	\end{frame}
	
	\begin{frame}{Q98 Theory --- Landscape}
		\begin{block}{Theory}
			\textbf{Landscape:}
			\begin{itemize}
				\item Relationships
				\item Critical models
				\item Data flows
				\item Interdependencies
			\end{itemize}
			\textbf{Why:}
			\begin{itemize}
				\item Cascading failure
				\item Correlated errors
				\item Systemic risk
			\end{itemize}
		\end{block}
	\end{frame}
	
	\begin{frame}{Q99 --- Reporting (repeat)}
		\begin{block}{Question}
			Purpose of model risk reporting?
		\end{block}
		\begin{enumerate}[label=\Alph*.]
			\item Increase accuracy
			\item Inform management about model risk
			\item Speed up training
			\item Reduce features
		\end{enumerate}
		\begin{exampleblock}{Answer: B}
		\end{exampleblock}
		\begin{block}{Explanation}
			Informs management about model risk.
		\end{block}
	\end{frame}
	
	\begin{frame}{Q99 Theory --- Reporting}
		\begin{block}{Theory}
			\textbf{Reporting:}
			\begin{itemize}
				\item Inventory summary
				\item Validation status
				\item Key issues
				\item Performance
				\item Remediation
			\end{itemize}
			\textbf{Audience:} Board, management, risk committee, regulators.
		\end{block}
	\end{frame}
	
	\begin{frame}{Q100 --- Change Management (repeat)}
		\begin{block}{Question}
			Purpose of model change management?
		\end{block}
		\begin{enumerate}[label=\Alph*.]
			\item Increase accuracy
			\item Control changes with review
			\item Speed up training
			\item Reduce features
		\end{enumerate}
		\begin{exampleblock}{Answer: B}
		\end{exampleblock}
		\begin{block}{Explanation}
			Controls changes with review, approval, documentation.
		\end{block}
	\end{frame}
	
	\begin{frame}{Q100 Theory --- Change Management}
		\begin{block}{Theory}
			\textbf{Change management:}
			\begin{itemize}
				\item Identify need
				\item Assess impact
				\item Approval
				\item Implement
				\item Document
				\item Validate
			\end{itemize}
			\textbf{Types:}
			\begin{itemize}
				\item Minor: parameters
				\item Major: methodology
				\item Emergency: urgent
			\end{itemize}
		\end{block}
	\end{frame}
	
	% ============================================================
	\section{Final Review}
	% ============================================================
	
	\begin{frame}{Key Takeaways}
		\begin{block}{Module 1: AI and Risk}
			\begin{itemize}
				\item Inscrutability vs explainability vs interpretability
				\item Classical AI vs deep learning
				\item Data types, scaling, PCA
			\end{itemize}
		\end{block}
		\begin{block}{Module 2: Tools and Techniques}
			\begin{itemize}
				\item K-means, hierarchical, DBSCAN
				\item WCSS, BCSS, silhouette, scree
				\item Random forest, SVM, NLP, regularization
				\item Metrics: precision, recall, F1, AUC
			\end{itemize}
		\end{block}
		\begin{block}{Module 3: Risk and Risk Factors}
			\begin{itemize}
				\item Sensitivity vs specificity
				\item Algorithmic bias sources
				\item Fairness metrics
			\end{itemize}
		\end{block}
	\end{frame}
	
	\begin{frame}{Key Takeaways (cont.)}
		\begin{block}{Module 4: Responsible and Ethical AI}
			\begin{itemize}
				\item Ethical frameworks: consequentialism, deontology, virtue
				\item Principles: beneficence, nonmaleficence, justice, autonomy, explainability
				\item Bias, fairness, privacy
			\end{itemize}
		\end{block}
		\begin{block}{Module 5: Data and AI Model Governance}
			\begin{itemize}
				\item Data governance, model governance
				\item Validation, monitoring, inventory, landscape
				\item Roles and responsibilities
			\end{itemize}
		\end{block}
	\end{frame}
	
	


	
\end{document}